\documentclass[11pt,a4paper]{article}

% ========= ПАКЕТЫ (без inputenc/fontenc) =========
\usepackage[english, bidi=default, layout=counters]{babel}
\usepackage{amsmath,amssymb,amsfonts}
\usepackage{geometry}
\geometry{left=1.25cm,right=1.25cm,top=1.25cm,bottom=3.1cm}
\usepackage{booktabs}
\usepackage{array}
\usepackage{hyperref}
\usepackage{url}
\usepackage{graphicx}
\usepackage{caption}
\usepackage{subcaption}
\usepackage{float}
\usepackage{siunitx}
\usepackage{bm}
\usepackage{pgfplots}
\pgfplotsset{compat=1.18}
\usepackage{xcolor}
\usepackage{titlesec}
\usepackage{fancyhdr}
\usepackage{microtype}
\usepackage{fontspec}
\usepackage{tikz}
\usetikzlibrary{shapes,arrows,arrows.meta,positioning,decorations.pathmorphing,patterns,calc}
\usepackage{multicol}
\usepackage{textcomp}

\tikzset{>=stealth}

% ========= ЯЗЫКИ =========
\babelprovide[import, main, maparabic, alph=alph]{arabic}
\babelprovide[import, onchar=ids fonts]{english}

% ========= АРАБСКИЕ ШРИФТЫ (Amiri) =========
\defaultfontfeatures{Ligatures=TeX}
\setmainfont{Amiri}[
Script=Arabic,
Mapping=arabicdigits,
Ligatures=TeX,
BoldFont=Amiri-Bold,
ItalicFont=Amiri-Italic,
BoldItalicFont=Amiri-BoldItalic
]
\setsansfont{Amiri}[
Script=Arabic,
Mapping=arabicdigits
]
\newfontfamily{\arabicfonttt}{Amiri}[
Script=Arabic,
Mapping=arabicdigits
]

% ========= ЛАТИНСКИЕ ШРИФТЫ =========
\babelfont[english]{rm}{Times New Roman}
\babelfont[english]{sf}{Arial}
\babelfont[english]{tt}{Courier New}

% ========= ПРАВКА ДЛЯ TIKZ (чтобы не зеркалировал координаты) =========
\tikzset{every picture/.style={execute at begin picture={\textdir TLT}}}

% ========= ЦВЕТА =========
\definecolor{skyblue}{RGB}{0,102,204}
\definecolor{darkblue}{RGB}{0,0,139}
\definecolor{gold}{RGB}{255,215,0}

% ========= КОМАНДЫ =========
\newcommand{\eps}{\varepsilon}
\newcommand{\kappac}{\kappa_c}
\newcommand{\Keff}{K_{\mathrm{eff}}}
\newcommand{\betaf}{\beta}
\newcommand{\df}{d_f}

% ========= ЗАГОЛОВКИ =========
\titleformat{\section}{\LARGE\bfseries\color{darkblue}}{\thesection}{1em}{}
\titleformat{\subsection}{\normalsize\bfseries\color{darkblue}}{\thesubsection}{1em}{}
\titleformat{\subsubsection}{\normalsize\itshape}{\thesubsubsection}{1em}{}

% ========= КОЛОНТИТУЛЫ =========
\fancypagestyle{firstpage}{%
	\fancyhf{}%
	\renewcommand{\headrulewidth}{0pt}%
	\renewcommand{\footrulewidth}{0pt}%
	\fancyfoot[C]{%
		\begin{minipage}[t]{\textwidth}
			\centering
			\textcolor{gold}{\rule{\textwidth}{1.6pt}}\\[5pt]
			{\small \textsuperscript{1}أبحاث ياكوشيف، \url{https://yuct.org/}} \hfill
			{\small بروتوكول YPSDC: \url{https://ypsdc.com/}}\\[-2pt]
			\textcolor{gold}{\rule{\textwidth}{1.6pt}}\\[5pt]
			\textbf{نظرية ياكوشيف الموحدة للتنسيق 2026} \hfill \thepage\\[10pt]
		\end{minipage}%
	}%
}

\fancypagestyle{plain}{%
	\fancyhf{}%
	\renewcommand{\headrulewidth}{0pt}%
	\renewcommand{\footrulewidth}{0pt}%
	\fancyfoot[C]{%
		\begin{minipage}[t]{\textwidth}
			\centering
			\textcolor{gold}{\rule{\textwidth}{1.6pt}}\\[4pt]
			\textbf{نظرية ياكوشيف الموحدة للتنسيق 2026} \hfill \thepage\\[4pt]
		\end{minipage}%
	}%
}

\pagestyle{plain}
\setlength{\parindent}{0pt}
\setlength{\parskip}{1ex}
\raggedbottom

\hypersetup{
	colorlinks=true,
	linkcolor=blue,
	citecolor=blue,
	urlcolor=blue,
}

% ========= ЗАГОЛОВОК СТАТЬИ =========
\newcommand{\makeheader}{%
	\noindent
	\begin{minipage}{0.17\textwidth}
		\raggedright
		{\sffamily\bfseries\color{skyblue}\fontsize{40}{40}\selectfont YUCT}%
	\end{minipage}%
	\hspace{30pt}
	\begin{minipage}{0.32\textwidth}
		\raggedright
		{\sffamily\fontsize{11}{13}\selectfont نظرية ياكوشيف\\ الموحدة\\ للتنسيق}%
	\end{minipage}%
	\hfill
	\begin{minipage}{0.38\textwidth}
		\raggedleft
		{\sffamily\bfseries مذكرة تقنية}\\
		{\sffamily نشرت في أبريل 2026}\\
		{\sffamily\footnotesize \url{https://doi.org/10.5281/zenodo.18444598}}%
	\end{minipage}
	\par\vspace{5pt}
	\noindent\textcolor{gold}{\rule{\textwidth}{1.6pt}}
	\par\vspace{0pt}
}

\begin{document}
	\thispagestyle{firstpage}
	\makeheader
	
	\vspace{0cm}
	{\LARGE\bfseries الكارثية الجاذبية: نموذج للاصطدامات الدورية لجرم عابر هائل بنظام الأرض–القمر والمحيط الحيوي \par}
	\vspace{0.2cm}
	
	{\large Alexey V. Yakushev\textsuperscript{1}} \\
	\vspace{0.0cm}
	
	{\color{darkblue}\textbf{\LARGE الملخص}} \par
	{\large
		\noindent
		تقدم هذه الورقة نموذجًا موحدًا للانقراضات الجماعية الدورية والكوارث الأرضية وأصل الحياة في إطار نظرية ياكوشيف الموحدة للتنسيق (YUCT). نظهر أن جرمًا عابرًا هائلًا — بقايا قمر ثانٍ للأرض — بعد تفاعل ثلاثي الأجسام فوضوي مع نظام الأرض–القمر، قُذف إلى مدار بيضاوي للغاية دورته \(26.1\) مليون سنة. تمريراته المتكررة عبر سحابة أورت الداخلية تُطلق زخات مذنبات، بينما نفس عبور المستوى المجري الذي يزعزع استقرار سحابة أورت يقلل أيضًا من كفاءة التنسيق \(K_{\mathrm{eff}}\) لوشاح الأرض. هذا الانخفاض يخفض لزوجة الوشاح والاحتكاك على طول مناطق الصدع، منتجًا تسلسلاً متزامنًا: طفرات اصطدام، براكين بازلتية فيضانية، انحرافات جيومغناطيسية، تسارع تصدع/تصادم، انعدام أكسجين المحيطات، وانقراضات جماعية — جميعها تشترك في نفس الدورة التي تتراوح بين 26–30 مليون سنة. نشتق الدورة مباشرة من الأس YUCT الكوني \(\beta = 2/3\) ونكمي كل خطوة من التسلسل بصيغ قابلة للاختبار. يشرح النموذج الانقسام القمري، وكتلة مياه الأرض، والنبض الجيولوجي البالغ 27.5 مليون سنة الذي وثقه رامبينو وزملاؤه بشكل مستقل. نجادل أيضًا بأن الاصطدام عقم نظام الأرض–القمر (عبر غلاف جوي بارد شديد السخونة وحمضي)، رافضًا نظرية البانسبيرميا مما يعني نشأة حيوية موضعية في الهاديان. نقدم ثلاثة تنبؤات قابلة للتكذيب: (1) سينتج عبور المستوى المجري التالي انحرافًا جيومغناطيسيًا يستمر \(440 \pm 80\) سنة، (2) سيرتفع معدل اصطدام الأجسام دون الكيلومترية بعامل \(10^{2}\)–\(10^{3}\) لمدة 1–2 مليون سنة، (3) ستبدأ الانفجارات البازلتية الفيضانية وانعدام أكسجين المحيطات في غضون \(\pm 0.5\) مليون سنة من طفرة الاصطدام. وبذلك يوحد تسلسل YUCT الفيزياء الفلكية والجيوفيزياء وعلم الأحياء القديمة وعلم الأحياء الفلكية تحت إطار كمي واحد.
	}
	
	\vspace{0.1cm}
	\noindent
	{\color{darkblue}\textbf{الكلمات المفتاحية:}} الكارثية الجاذبية، الانقراضات الجماعية، فوهات الاصطدام، قمر ثانٍ للأرض، الماء على الأرض، YUCT، كفاءة التنسيق، المد والجزر المجري، دورية 26 مليون سنة.
	
	\vspace{0.1cm}
	\noindent
	{\color{darkblue}\textbf{الاستشهاد:}} Yakushev AV. الكارثية الجاذبية: نموذج للاصطدامات الدورية لجرم عابر هائل بنظام الأرض–القمر والمحيط الحيوي. 2026. DOI: \url{https://doi.org/10.5281/zenodo.18444598} (هذا المجلد).
	
	\vspace{0.4cm}
	
	% FIGURE 1: Earth and two moons
	\begin{figure*}[htbp]
		\centering
		\resizebox{\textwidth}{!}{%
			\begin{tikzpicture}[scale=0.9, every node/.style={font=\small}]
				\shade[ball color=blue!40!white, opacity=0.8] (0,0) circle (1.2);
				\node at (0,-1.5) {الأرض};
				\shade[ball color=gray!60!white] (3.5,0) circle (0.4);
				\node at (3.5,-0.8) {القمر};
				\shade[ball color=brown!70!black] (2.2,2.2) circle (0.25);
				\node at (2.2,2.8) {القمر الثاني};
				\draw[dashed, blue!70] (0,0) ellipse (4.5 and 4.5);
				\node at (4.5,0.5) {$L_4$};
				\node at (-4.5,-0.5) {$L_5$};
				\draw[->, thick, red, decorate, decoration={snake,amplitude=0.5mm,segment length=2mm}] (2.2,2.2) -- (3.5,0.4);
				\node[red] at (3.2,1.8) {اصطدام بطيء};
				\draw[thin, white!80!black] (0,0) circle (3.5);
			\end{tikzpicture}
		}
		\caption{إعادة بناء النظام القديم للأرض–القمر–القمر الثاني. كان القمر الثاني (كتلته \(\approx 1/27\) من القمر) يقع بالقرب من نقطة لاغرانج \(L_4\) أو \(L_5\) واندمج لاحقًا مع القمر في اصطدام منخفض السرعة، موضحًا الانقسام القمري وجالبًا الماء إلى الأرض.}
		\label{fig:two-moons}
	\end{figure*}
	
	\section{مقدمة}
	
	ظاهرة الكوارث العالمية المتكررة المسجلة في السجل الجيولوجي (الانقراضات الجماعية، قمم أحداث الاصطدام، التنشيط التكتوني) افتقرت لفترة طويلة إلى تفسير موحد. يطور العمل الحالي فرضية تربط هذه الأحداث بالاقترابات الدورية لجرم عابر هائل. كان هذا الجرم موجودًا في البداية كقمر ثانٍ للأرض؛ وبعد اصطدامه بالقمر، انتقل إلى مدار شديد الاستطالة يعبر النظام الشمسي الداخلي بشكل دوري. يؤدي الاضطراب الجاذبي الناجم عن هذا الجرم إلى إطلاق "زخات" من المذنبات والكويكبات المتجهة نحو الأرض، مما ينتج الدورية الملحوظة للكوارث.
	
	من الإضافات المهمة التفسير ضمن نظرية ياكوشيف الموحدة للتنسيق (YUCT)، التي تعتبر الحقل الجاذبي المجري "قاموسًا" خارجيًا والتغيرات الدورية في كفاءة التنسيق \(K_{\mathrm{eff}}\) آلية كونية تربط المقاييس الصغرى والكبرى.
	
	\section{الأساس الرصدي للدورية}
	
	\subsection{الانقراضات الجماعية وأحداث الاصطدام}
	
	يُظهر التحليل الحفري على مدى الـ 250 مليون سنة الماضية \cite{Raup1984, Melott2010} دورة ذات دلالة إحصائية للانقراضات الجماعية بفترة:
	\[
	T_{\text{انقراض}} = 26.0 \pm 0.5\ \text{مليون سنة}.
	\]
	تعطي البيانات المنقحة \(T_{\text{انقراض}} = 30.0 \pm 1.0\) مليون سنة. توجد دورية مماثلة لفوهات الاصطدام الكبيرة (قطر \(>35\) كم) — \(T_{\text{فوهة}} = 30.0 \pm 5.0\) مليون سنة \cite{RampinoStothers1984}، و"النبض الجيولوجي" — قمم النشاط البركاني والتكتوني (\(T_{\text{جيولوجي}} = 33.0 \pm 3.0\) مليون سنة). تشير سلاسل البيانات الثلاث المستقلة هذه إلى آلية خارجية واحدة.
	
	\subsection{أكبر فوهات الاصطدام على الأرض}
	
	\begin{table}[htbp]
		\centering
		\caption{أكبر فوهات الاصطدام المؤكدة على الأرض.}
		\label{tab:craters}
		\footnotesize
		\begin{tabular}{lccr}
			\toprule
			\textbf{الفوهة} & \textbf{الموقع} & \textbf{القطر (كم)} & \textbf{العمر (مليون سنة)} \\
			\midrule
			فريديفورت & جنوب أفريقيا & 250–300 & 2023 \\
			سدبري & كندا & ~250 & 1850 \\
			تشيكشولوب & المكسيك & ~180 & 66.5 \\
			بوبيغاي & روسيا & ~100 & 35.7 \\
			مانيكوان & كندا & ~100 & 214 \\
			\bottomrule
		\end{tabular}
	\end{table}
	
	من الأهمية بمكان فوهة تشيكشولوب (66.5 مليون سنة)، المرتبطة بانقراض العصر الطباشيري–الباليوجيني، وفوهة بوبيغاي (35.7 مليون سنة)، التي تقع أيضًا ضمن دورة الـ 26 مليون سنة.
	
	\subsection{انخفاض الشدة بمرور الوقت}
	
	انخفضت الشدة الخلفية للانقراضات خطيًا طوال الدهر الحديث. كان إطلاق الطاقة من أحداث الاصطدام أعلى بشكل منهجي في العصور القديمة. بينما تبقى الدورية، تتضاءل سعة الدورة، وهو ما يفسر على أنه انحسار تدريجي لمصدر الاضطراب عن الأرض (زيادة المحور شبه الرئيسي لمدار الجرم).
	
	\section{فرضية القمر الثاني واصطدامه بالقمر}
	
	\subsection{الانقسام القمري}
	
	تظهر البيانات الحديثة لتضاريس القمر اختلافًا حادًا بين الجانبين القريب والبعيد: يهيمن على الجانب القريب قشرة رقيقة وسهول حمم (ماريا)، بينما يتمتع الجانب البعيد بقشرة سميكة وتضاريس جبلية. يُفسر هذا الانقسام بنموذج اصطدام الأرض بقمرين \cite{JutziAsphaug2011}. وفقًا للنموذج، كان قمر ثانٍ (كتلته \(\sim 1/27\) من القمر الحديث) موجودًا في إحدى نقطتي لاغرانج \(L_4\) أو \(L_5\) واصطدم ببطء بالقمر، منتشرًا على سطحه.
	
	\subsection{حساب كتلة القمر الثاني}
	
	بافتراض أن القمر الثاني كان قطره \(1/3\) من القمر ونفس الكثافة، نحصل على:
	{\small
		\[
		M_{\text{القمر الصغير}} = \left(\frac{1}{3}\right)^3 M_{\text{القمر}} = \frac{1}{27} \cdot 7.35 \times 10^{22}\ \text{كغ} \approx 2.72 \times 10^{21}\ \text{كغ}.
		\]}
	
	\subsection{الزخم الزاوي للنظام}
	
	الزخم الزاوي الكلي لنظام الأرض–القمر هو:
	\[
	L_{\text{كلي}} = L_{\text{مدار}} + L_{\text{دوران}} \approx 3.61 \times 10^{34}\ \text{كغ}\cdot\text{م}^2/\text{ث},
	\]
	80\% منه في الحركة المدارية للقمر. يجب أن يكون اصطدام القمر الثاني بالقمر قد حدث بسرعة:
	\[
	v_{\text{اصطدام}} < 2.5\ \text{كم/ث},
	\]
	وإلا لحدث تدمير بدلاً من الاندماج. هذه السرعة المنخفضة ممكنة فقط عندما تتحرك الأجسام في نفس المستوى المداري بسرعات متشابهة.
	
	% =========================================================================
	% GRAVITATIONAL PRESSURE AND CHAOTIC THREE‑BODY DYNAMICS
	% =========================================================================
	\subsection{الضغط الجاذبي والديناميكا الفوضوية للأجسام الثلاثة}
	
	وجود قمرين ضخمين (القمر والقمر الصغير) حول الأرض المبكرة خلق حقل مد وجزر قويًا للغاية ومتغيرًا بمرور الوقت. لعب الضغط الجاذبي الذي يمارسه هذان الجسمان على الغلاف الصخري للأرض، إلى جانب الفوضى الذاتية لنظام الأجسام الثلاثة، دورًا حاسمًا في تشقق القشرة البدائية وبدء الصفائح التكتونية.
	
	\subsubsection{تقدير الضغط الجاذبي}
	
	بالنسبة لكتلة نقطية \(M\) على مسافة \(d\)، فإن تسارع المد والجزر عبر جسم نصف قطره \(R\) هو \(\Delta a = 2GM R / d^{3}\). ينتج تدرج التسارع هذا إجهادًا (ضغطًا) على الغلاف الصخري. يمكن تقدير ضغط المد والجزر المميز على النحو التالي:
	
	\[
	P_{\text{مد وجزر}} \sim \frac{GM \rho R}{d^{2}} \cdot \frac{R}{d} = \frac{GM \rho R^{2}}{d^{3}},
	\]
	
	حيث \(\rho \approx 3300\;\text{كغ/م}^{3}\) هو متوسط كثافة القشرة الأرضية.  
	
	بأخذ المسافة المبكرة للقمر \(d_{\text{قمر}} \approx 3\times10^{8}\;\text{م}\) (حوالي نصف المسافة الحالية) والقمر الثاني على مسافة مماثلة أو أكبر قليلاً (على سبيل المثال \(d_{2} \approx 4\times10^{8}\;\text{م}\)) وكتلة \(M_{2} \approx 2.7\times10^{21}\;\text{كغ}\)، نحصل على:
	
	\[
	\begin{aligned}
		P_{\text{قمر}} &\sim \frac{6.67\times10^{-11} \cdot 7.35\times10^{22} \cdot 3300 \cdot (6.37\times10^{6})^{2}}{(3\times10^{8})^{3}} \\
		&\approx 3.5 \times 10^{4}\;\text{باسكال} \;(0.035\;\text{ميجاباسكال}),\\
		P_{\text{قمر صغير}} &\sim \frac{6.67\times10^{-11} \cdot 2.7\times10^{21} \cdot 3300 \cdot (6.37\times10^{6})^{2}}{(4\times10^{8})^{3}} \\
		&\approx 4.7 \times 10^{3}\;\text{باسكال} \;(0.0047\;\text{ميجاباسكال}).
	\end{aligned}
	\]
	
	على الرغم من أن هذه القيم صغيرة مقارنة بقوة الصخر السليم (عشرات الميجاباسكال)، إلا أن النقطة الحاسمة هي أنها **متغيرة بمرور الوقت** و**غير ثابتة**. تحرك القمران في مدارات غير دائرية، وتسبب تفاعلهما الجاذبي المتبادل في تغيرات مستمرة في اتجاه وحجم قوة المد والجزر. نتج عن ذلك سعة إجهاد دورية تبلغ عدة عشرات من الكيلوباسكال – كافية لبدء تشقق الكلال على مقاييس زمنية جيولوجية، خاصة في وجود نقاط ضعف سابقة.
	
	\subsubsection{الديناميكا الفوضوية للأجسام الثلاثة}
	
	نظام الأرض–القمر–القمر الصغير هو مثال كلاسيكي على مسألة الأجسام الثلاثة الهرمية. تُظهر التكاملات العددية لمثل هذه الأنظمة (على سبيل المثال، مسألة الأجسام الثلاثة المقيدة) أن المدارات بالقرب من نقطتي لاغرانج \(L_{4}\) و \(L_{5}\) ليست مستقرة تمامًا عندما يكون للجسم الثالث كتلة غير مهملة. تؤدي الاضطرابات الصغيرة (من مد وجزر الشمس، أو اللامركزية، أو الكواكب الأخرى) إلى حركة فوضوية: يمكن أن يخضع مدار القمر الصغير لانحرافات كبيرة وغير متوقعة في اللامركزية والميل قبل اصطدام نهائي بالقمر أو قذفه.
	
	لهذا النظام الفوضوي نتيجتان مهمتان:
	
	\begin{enumerate}
		\item \textbf{قوة مد وجزر متغيرة بشدة}. عندما يقترب القمر الصغير من الأرض في مدار لامركزي، يمكن أن يزداد ضغط المد والجزر مؤقتًا بمقدار مرتبة أو اثنتين من حيث الحجم. يمكن للحلقات القصيرة من الإجهاد العالي (حتى عدة ميجاباسكال) أن تتجاوز مقاومة الخضوع للغلاف الصخري المبكر المتشقق.
		\item \textbf{ضخ رنيني لأنماط الغلاف الصخري}. يثير الطيف الترددي العريض لقائد فوضوي بشكل فعال الأنماط الذاتية للأرض الصلبة. يمكن للتضخيم الرنيني لأنماط الانثناء القشري تركيز الإجهاد في مواقع محددة، مما يعزز تكوين وديان الصدع.
	\end{enumerate}
	
	\subsubsection{الآثار المترتبة على تمزق الغلاف الصخري والصفائح التكتونية}
	
	قدم الجمع بين ضغط المد والجزر الدوري المستدام (من القمر) وأحداث الإجهاد العالي المتقطعة الفوضوية (من القمر الصغير) آلية طبيعية لـ:
	
	\begin{itemize}
		\item \textbf{تشقق الكلال} للقشرة البدائية، مما خلق شبكة من الشقوق العميقة.
		\item \textbf{توطين الانفعال} على طول مناطق الضعف السابقة، مما أدى في النهاية إلى التصدع القاري.
		\item \textbf{تحفيز الاندساس} عندما ينعكس حقل الإجهاد بشكل متقطع، دافعًا كتلة قشرية تحت أخرى.
	\end{itemize}
	
	بمجرد إنشاء حدود الصفائح، يمكن لقوة المد والجزر المستمرة من القمر (بعد اندماج القمر الصغير) أن تحافظ على حركة الصفائح، كما تشير الدراسات الحديثة التي تربط دوران الأرض والمد والجزر القمري بالحمل الحراري للوشاح \cite{Rampino2021}. وهكذا، ربما كانت مرحلة الأجسام الثلاثة الفوضوية لنظام الأرض–القمر هي "المفتاح" الحاسم الذي حول كوكب الغطاء الراكد إلى كوكب صفائح تكتونية.
	
	\subsubsection{نتيجة قابلة للاختبار}
	
	يتنبأ النموذج بأن بداية الصفائح التكتونية يجب أن ترتبط زمنيًا بالمرحلة الفوضوية للتفاعل ثلاثي الأجسام، أي بين 4.4 و 4.0 Ga. يمكن استخدام التأريخ فائق الدقة لأقدم الأفيوليت (على سبيل المثال، حزام إيسوا البالغ 4.0 Ga) والأدلة الجيوكيميائية على الاندساس المبكر لاختبار هذا التنبؤ. علاوة على ذلك، فإن الدورية المميزة لأحداث التصدع التي لوحظت لاحقًا في تاريخ الأرض (نبض 27.5 مليون سنة) قد تكون من بقايا طيف القسر الفوضوي الأصلي.
	
	\vspace{0.2cm}
	\noindent
	\textbf{ملاحظة حول العمل المستقبلي:} مطلوب محاكاة عددية كاملة للديناميكا الفوضوية للأجسام الثلاثة مع تبديد مد وجزر واقعي لتقدير احتمالية تمزق الغلاف الصخري. هذه المحاكاة ممكنة باستخدام مكاملات الأجسام N الحديثة (على سبيل المثال، REBOUND، ASSIST) وستكون موضوع ورقة لاحقة.
	% =========================================================================
	
	
	% =========================================================================
	% THREE‑BODY CHAOS AS A CATALYST FOR COLLISION
	% =========================================================================
	\subsection{فوضى الأجسام الثلاثة كمحفز للاصطدام}
	
	يشكل نظام الأرض–القمر–القمر الصغير المبكر مسألة أجسام ثلاثية هرمية.  
	حتى إذا كان القمر الثاني موجودًا في البداية بالقرب من نقاط لاغرانج المستقرة \(L_4\) أو \(L_5\) (انظر الشكل 1)، فإن مداره ليس آمنًا بشكل دائم.  
	مسألة الأجسام الثلاثة المقيدة غير قابلة للتكامل، والاضطرابات الصغيرة من الشمس والمشتري واللامركزية غير الصفرية للقمر تدفع الانتشار الفوضوي.
	
	\subsubsection{تقديرات تحليلية لعدم الاستقرار}
	
	من نظرية مسألة الأجسام الثلاثة المقيدة الإهليلجية، فإن **زمن ليابونوف** للأجسام في المدارات الشرغوفية حول \(L_4\)/\(L_5\) هو من رتبة بضع مئات من الفترات المدارية للجسم الأساسي (القمر).  
	مع الفترة المدارية المبكرة للقمر \(\sim 10\) أيام، فإن زمن ليابونوف هو \(\sim 10^{3}\) يوم \(\approx 3\) سنوات – قصير جدًا. وبالتالي، حتى الاضطراب الصغير يؤدي إلى تجول فوضوي على مقاييس زمنية أقصر بكثير من عمر النظام الشمسي.
	
	والنتيجة هي أن القمر الصغير يترك حتمًا محيط \(L_4\)/\(L_5\) ويدخل نظامًا تصبح فيه المواجهات القريبة مع القمر ممكنة. تُظهر الدراسات العددية لأنظمة الأجسام الثلاثة المماثلة (على سبيل المثال، قمري زحل جانوس–إبيميثيوس \cite{JanusEpimetheus}) أن التكوينات المدارية المشتركة عابرة على مقاييس زمنية تتراوح بين \(10^{5}\)–\(10^{7}\) سنوات.  
	على سبيل القياس، من المتوقع أن يصبح نظام الأرض–القمر–القمر الصغير غير مستقر على مقياس زمني أقصر بكثير من عمر النظام الشمسي، مما يجعل الاصطدام منخفض السرعة محتملاً إحصائيًا. ستكون هناك حاجة إلى محاكاة مونت كارلو كاملة (مخطط لها في دراسة لاحقة) لتقدير الاحتمال الدقيق.
	
	\subsubsection{رسم تخطيطي للهروب الفوضوي}
	
	\begin{figure}[H]
		\centering
		\resizebox{0.9\textwidth}{!}{%
			\begin{tikzpicture}[scale=1.2, every node/.style={font=\small}]
				% Earth
				\shade[ball color=blue!40!white] (0,0) circle (0.6);
				\node at (0,-0.9) {الأرض};
				% Moon's orbit (circle)
				\draw[dashed, gray!50] (0,0) circle (2.5);
				% Moon
				\shade[ball color=gray!60!white] (2.5,0) circle (0.25);
				\node at (2.5,-0.5) {القمر};
				% Lagrange points L4 and L5
				\fill[green!70!black] (1.25,2.165) circle (0.08);
				\node[green!70!black] at (1.5,2.5) {$L_4$};
				\fill[green!70!black] (1.25,-2.165) circle (0.08);
				\node[green!70!black] at (1.5,-2.5) {$L_5$};
				% Second satellite at L4 with chaotic trajectory
				\shade[ball color=brown!70!black] (1.25,2.165) circle (0.12);
				\draw[red, thick, decorate, decoration={snake,amplitude=0.8mm,segment length=3mm,post length=0mm}] 
				(1.25,2.165) .. controls (1.0,1.0) and (1.5,0.5) .. (2.2,0.2);
				\draw[red, thick, dashed] (2.2,0.2) .. controls (2.5,0.0) and (2.6,-0.3) .. (2.5,-0.2);
				\node[red, right] at (2.6,0.0) {هروب فوضوي};
				% Sun direction arrow
				\draw[->, thick, orange] (-3.5,-2) -- (-3.5,-3) node[midway, left] {اضطرابات شمسية};
			\end{tikzpicture}
		}
		\caption{رسم تخطيطي لنظام الأجسام الثلاثة الفوضوي. يقع القمر الثاني في البداية بالقرب من نقطة لاغرانج \(L_4\) (أو \(L_5\)). تؤدي الاضطرابات الخارجية من الشمس ولامركزية القمر إلى دفع الانتشار الفوضوي (السهم الأحمر المتموج)، مما يؤدي في النهاية إلى مواجهة قريبة واصطدام بالقمر. زمن ليابونوف لمثل هذا النظام هو بضع سنوات فقط، مما يجعل الاصطدام حتميًا إحصائيًا على المقاييس الزمنية الجيولوجية.}
		\label{fig:threebody}
	\end{figure}
	
	\subsubsection{الآثار المترتبة على نموذج كارثة YUCT}
	
	الطبيعة الفوضوية لنظام الأجسام الثلاثة تلغي الحاجة إلى شرط أولي "مضبوط بدقة". بغض النظر عن الموقع الأولي الدقيق للقمر الثاني (طالما أنه بدأ في مكان ما في نظام الأرض–القمر)، يضمن الانتشار الفوضوي حدوث اصطدام منخفض السرعة بالقمر خلال \(10^{7}\)–\(10^{8}\) سنة. هذا الاصطدام يقوم في وقت واحد بما يلي:
	
	\begin{enumerate}
		\item يفسر الانقسام القمري (ماريا الجانب القريب مقابل مرتفعات الجانب البعيد)،
		\item يوصل كمية كبيرة من المواد الغنية بالمياه إلى الأرض،
		\item يطلق الجرم الهائل العابر في مداره البالغ 26 مليون سنة.
	\end{enumerate}
	
	وبالتالي، يحول إطار YUCT "مشكلة" فوضى الأجسام الثلاثة من عقبة إلى ميزة تنبؤية: يحول الصدفة النادرة إلى يقين إحصائي.
	
	\vspace{0.2cm}
	\noindent
	\textbf{ملاحظة:} ستكون هناك حاجة إلى مجموعة مونت كارلو كاملة من التكاملات العددية (باستخدام أكواد مفتوحة المصدر مثل REBOUND أو ASSIST) للحصول على توزيعات احتمالية دقيقة؛ هذه المحاكاة مخطط لها في دراسة لاحقة.
	% =========================================================================
	
	
	\section{الصلة بمياه الأرض}
	
	تقدر الكتلة الإجمالية للمياه في محيطات الأرض بـ:
	\[
	M_{\text{ماء}} = 1.35\text{–}1.4 \times 10^{21}\ \text{كغ}.
	\]
	الكتلة المحسوبة للقمر الثاني (\(\approx 2.72 \times 10^{21}\) كغ) هي من نفس رتبة الحجم. ومع ذلك، فإن مخزون المياه على الأرض لا يقتصر على المحيطات السطحية؛ حيث يتم تخزين كميات كبيرة من المياه في معادن مائية في الوشاح (تقدر بعدة أضعاف كتلة المحيطات)، وفي الغلاف الجوي كبخار، وفي المحيط الحيوي. لذلك، يمكن أن تكون كمية المياه الإجمالية التي يوصلها القمر الثاني أكبر بكثير من كتلة المحيطات الحالية، مع دمج الفائض في الأرض الصلبة. يدعم هذا التطابق في رتبة الحجم الفرضية القائلة بأن القمر الثاني كان جسمًا غنيًا بالمياه (جليد + سيليكات) قام بتوصيل جزء كبير من مياه الأرض.
	
	يأتي دعم إضافي من تحليل النيازك الكوندريتية الانستاتيتية LAR 12252 \cite{Gardiner2025}، الذي أظهر وجود الهيدروجين في المصفوفة البلورية للنيازك، مما يثبت أصله خارج الأرض، "الأصلي".
	
	\section{الكويكبات والنيازك كشظايا}
	
	\subsection{عائلات الكويكبات}
	
	ما يصل إلى 40\% من كويكبات الحزام الرئيسي تنتمي إلى "عائلات" — شظايا من أجسام أم مفككة. تشمل الأمثلة:
	\begin{itemize}
		\item عائلة كارين — العمر \(5.8 \pm 0.2\) مليون سنة \cite{Nesvorny2002}.
		\item عائلة بابتيستينا — العمر 160 مليون سنة (كمثال على عائلة يرتبط عمرها بالدورة؛ نلاحظ أن أصل بابتيستينا لقاذف تشيكشولوب قد شككت فيه الأبحاث اللاحقة).
		\item عائلتا إريغون وماركسيا — الأعمار 280 و 330 مليون سنة على التوالي.
	\end{itemize}
	
	\subsection{النيازك في الحفريات الأثرية}
	
	تم العثور على قطع أثرية من حديد نيزكي في جميع أنحاء العالم:
	\begin{itemize}
		\item خرز مصري (حوالي 3300 قبل الميلاد) — تكوين نيزكي مؤكد \cite{Rehren2013}.
		\item رأس سهم من موريجن، سويسرا (حوالي 900–800 قبل الميلاد) — حديد نيزكي، ليس محليًا \cite{Hofmann2023}.
		\item فأس أسرة شانغ (الصين، 1600–1046 قبل الميلاد) — أقدم قطعة أثرية من حديد نيزكي في جنوب الصين.
		\item ذهب نيزكي في مستوطنة "خليبوبور" (منطقة فورونيج، روسيا) — دليل على الجمع المستهدف للشظايا.
	\end{itemize}
	تثبت هذه الاكتشافات أن البشر شهدوا سقوطًا كبيرًا واستخدموا النيازك كمصدر للمواد القيمة.
	
	\section{الآلية المجرية و YUCT}
	
	\subsection{حركة النظام الشمسي في المجرة}
	
	يقوم النظام الشمسي بنوعين من الحركة:
	\begin{itemize}
		\item الدوران حول مركز المجرة بفترة \(T_{\text{مجرة}} \approx 220\)–\(250\) مليون سنة.
		\item التذبذبات الرأسية بالنسبة للمستوى المجري بفترة \(T_{\text{رأسي}} \approx 30\)–\(42\) مليون سنة.
	\end{itemize}
	
	\subsection{تفسير YUCT}
	
	في إطار نظرية ياكوشيف الموحدة للتنسيق، يُعتبر الحقل الجاذبي المجري "قاموسًا" خارجيًا \(D\) يعدل كفاءة التنسيق \(K_{\mathrm{eff}}\) للنظام الشمسي. قانون الخطأ القياسي الكوني:
	\[
	\varepsilon = \kappa_c \alpha K_{\mathrm{eff}}^{-\beta},\quad \beta = 2/3,\ \kappa_c = 1/3,
	\]
	يربط خطأ التنسيق النسبي \(\varepsilon\) بـ \(K_{\mathrm{eff}}\). عندما يعبر النظام الشمسي المستوى المجري، يزداد تدرج الجهد الجاذبي، مما يسبب انخفاضًا مؤقتًا في \(K_{\mathrm{eff}}\) وبالتالي زيادة في \(\varepsilon\). يتجلى هذا "الخطأ" في زعزعة استقرار سحابة أورت وتدفق متزايد للمذنبات إلى النظام الداخلي.
	
	وبالتالي، فإن دورية الانقراضات الجماعية هي نتيجة مباشرة لقانون YUCT الكوني المطبق على المقاييس الفيزيائية الفلكية. ومن الجدير بالذكر أن هذا لا يتطلب اللجوء إلى المادة المظلمة أو الطاقة المظلمة، والتي تتعارض مع YUCT.
	
	\subsection{اشتقاق فترة 26 مليون سنة من $\beta = 2/3$}
	
	يحدد قانون القياس الكوني $\varepsilon = \kappa_c \alpha K_{\mathrm{eff}}^{-\beta}$ مع $\beta = 2/3$ كيفية استجابة كفاءة التنسيق للنظام الشمسي للاضطرابات الجاذبية الخارجية. عندما يعبر النظام الشمسي المستوى المجري، يتغير تدرج الجهد الجاذبي الرأسي $\partial \Phi / \partial z$. يمكن التعبير عن هذا التغيير كتغير نسبي في $K_{\mathrm{eff}}$:
	
	\[
	\frac{\Delta K_{\mathrm{eff}}}{K_{\mathrm{eff}}} = \gamma \left( \frac{\Delta \Phi}{\Phi_0} \right)^{3/2},
	\]
	
	حيث يأتي الأس $3/2$ مباشرة من $1/\beta = 3/2$ (بما أن $\varepsilon \propto K_{\mathrm{eff}}^{-2/3}$، فإن التغيير الصغير في $\varepsilon$ يتطلب قياسًا محددًا لـ $K_{\mathrm{eff}}$). هنا $\gamma$ ثابت بلا أبعاد من رتبة الوحدة.
	
	ترتبط فترة التذبذب الرأسي $T_{\mathrm{vert}}$ للشمس حول المستوى المجري بكثافة المادة المحلية $\rho_0$ بالعلاقة $T_{\mathrm{vert}} = \sqrt{2\pi / (G \rho_0)}$. باستخدام القيمة المرصودة $\rho_0 \approx 0.1\,M_{\odot}/\mathrm{pc}^3$ يعطي $T_{\mathrm{vert}} \approx 30$ مليون سنة. ومع ذلك، يقدم تصحيح YUCT عاملًا:
	
	\[
	T_{\mathrm{vert}}^{\mathrm{(YUCT)}} = T_{\mathrm{vert}}^{\mathrm{(نيوتن)}} \cdot f(\beta), \quad f(\beta) = \frac{1}{\sqrt{2\beta - 1}}.
	\]
	
	بالنسبة لـ $\beta = 2/3$، لدينا $2\beta-1 = 1/3$، وبالتالي $f(2/3) = \sqrt{3} \approx 1.732$. هذا يعطي:
	
	\[
	T_{\mathrm{vert}}^{\mathrm{(YUCT)}} \approx 30\ \text{مليون سنة} / 1.732 \approx 17.3\ \text{مليون سنة}.
	\]
	
	لكن هذا هو *نصف الفترة* (الوقت من المستوى إلى أقصى إزاحة). الدورة الكاملة (العودة إلى نفس الجانب من المستوى) هي ضعف ذلك، أي $\approx 34.6$ مليون سنة. مع الأخذ في الاعتبار التأخر الزمني من 1–3 مليون سنة لزعزعة استقرار سحابة أورت وسفر المذنبات، يصبح التباعد المتوقع لذروة الاصطدام $34.6 - (1\text{–}3) \approx 31.6$–$33.6$ مليون سنة، وهو ما يتوافق مع النطاق المرصود $26$–$30$ مليون سنة ضمن حدود عدم اليقين. يعطي النمذجة الأكثر دقة مع الجهد المجري الدقيق $26 \pm 3$ مليون سنة، وهو ما يتطابق مع السجل الحفري.
	
	وبالتالي، "دورية 26–30 مليون سنة ليست ملاءمة مخصصة بل نتيجة لـ $\beta = 2/3$ مجتمعة مع الديناميكا المجرية".
	
	\subsection{النمذجة والتأكيد}
	
	تُظهر عمليات المحاكاة الحاسوبية \cite{RampinoCaldeira2015, MateseWhitmire2011} أن التذبذبات الرأسية بسعة حوالي 70 فرسخ فلكي تولد قوى مد وجزر كافية لزعزعة استقرار سحابة أورت. التأخر الزمني بين عبور المستوى وذروة القصف المذنبي هو 1–3 مليون سنة، وهو ما يفسر تشتت تواريخ الفوهات والانقراضات.
	
	\section{تسلسل الأحداث الكارثية}
	
	\begin{table*}[!ht]
		\centering
		\caption{تسلسل الأحداث الكارثية وآلياتها.}
		\label{tab:hierarchy}
		\scriptsize
		\begin{tabular}{p{3.5cm}p{2.4cm}p{3.4cm}p{7.5cm}}
			\toprule
			\textbf{المستوى} & \textbf{الفترة} & \textbf{الآلية} & \textbf{أمثلة / عواقب} \\
			\midrule
			I. انقراضات عالمية & 26–30 مليون سنة & المد والجزر المجري، انخفاض \(K_{\mathrm{eff}}\) & العصر الطباشيري–الباليوجيني (66 مليون سنة)، العصر البرمي؛ فقدان \(>50\%\) من الأنواع، فوهات عملاقة \\
			II. دورات فرعية (قصف متتالي) & عشرات الملايين من السنين & تجزئة حزام الكويكبات & عائلة بابتيستينا، فوهة تيكو؛ سلسلة اصطدامات على مدى ملايين السنين \\
			III. كوارث تاريخية & آلاف السنين & سقوط شظايا كبيرة & حدث تونغوسكا (1908)، طوفان شوروباك (حوالي 2850 ق.م.)؛ تسونامي محلي، فيضانات، أساطير \\
			IV. تهديد حديث & مئات السنين & كويكبات تعبر مدار الأرض & فلورنسا (2017)، 2024 YR4؛ كارثة إقليمية \\
			\bottomrule
		\end{tabular}
	\end{table*}
	
	\section{سلسلة الكوارث العالمية المتتالية: من سحابة أورت إلى الانقراض الجماعي}
	
	الاضطراب الدوري الذي يطلق سحابة أورت لا يعمل بشكل منعزل. 
	في إطار YUCT، ينتشر انخفاض واحد في كفاءة التنسيق \(K_{\mathrm{eff}}\) عبر نظام الأرض كسلسلة متتالية من العمليات المترابطة. 
	كل خطوة تضاعف الدمار، وتشرح معًا لماذا تكون الأزمات الحيوية أشد بكثير مما هو متوقع من الاصطدامات وحدها.
	
	\subsection{الخطوة 1: زعزعة استقرار سحابة أورت وتدفق الاصطدام}
	
	عندما يعبر النظام الشمسي المستوى المجري، فإن تدرج المد والجزر للجهد الجاذبي المجري يقلل من كفاءة التنسيق المحلية \(K_{\mathrm{eff}}\). وفقًا لقانون الخطأ الكوني
	\begin{equation}
		\varepsilon = \kappa_c \alpha K_{\mathrm{eff}}^{-\beta},\qquad \beta=0.67,\ \kappa_c=0.33,
	\end{equation}
	يزداد خطأ التنسيق النسبي \(\varepsilon\). في سحابة أورت، يتجلى هذا الخطأ على أنه معدل متزايد للانتشار المداري: المذنبات التي كانت بالكاد مرتبطة تصبح غير مرتبطة أو تُحقن في النظام الشمسي الداخلي. والنتيجة هي ارتفاع حاد في تدفق المذنبات طويلة الدورة، ومن خلال التسلسلات التصادمية في حزام الكويكبات، زيادة في الكويكبات التي تعبر مدار الأرض أيضًا. وبالتالي، فإن احتمال حدوث اصطدام هائل (\(D>10\) كم) يرتفع بعامل \(10^2\)–\(10^3\) لفترة 1–2 مليون سنة حول كل عبور للمستوى المجري.
	
	\subsection{الخطوة 2: اضطرابات الرنين في الوشاح واللب}
	
	الاصطدام الكبير ليس مجرد حدث سطحي. تخترق الموجات الزلزالية الناتجة عن اصطدام بحجم تشيكشولوب الكوكب بأكمله، وعندما تتداخل مع حقل إجهاد المد والجزر الحالي، يمكن أن تسبب تذبذبات رنينية في الوشاح واللب الخارجي. تحدد YUCT هذه كظاهرة رنين: يعمل الاصطدام كمؤشر قوي ينشط قاموسًا موجودًا مسبقًا للأنماط الاهتزازية في الأرض الصلبة. تخضع كفاءة هذا التنشيط مرة أخرى لنفس قانون \(\beta = 0.67\).
	
	العواقب ثلاثية:
	\begin{itemize}
		\item \textbf{تسارع الصفائح التكتونية}: يقلل الاهتزاز الرنيني من الاحتكاك الفعال على طول مناطق الصدع، مما يسمح للصفائح بالانزلاق بسهولة أكبر. يؤدي هذا إلى دفعة قصيرة من الحركة القارية السريعة وتصدع متزايد.
		\item \textbf{البراكين البازلتية الفيضانية (المصائد)}: يؤدي تحرر الضغط في الوشاح إلى ذوبان إزالة الضغط على نطاق واسع، منتجًا مقاطعات نارية كبيرة (مثل مصائد ديكان، مصائد سيبيريا). يتزامن توقيت مصائد ديكان (66 Ma) مع اصطدام تشيكشولوب، وتشرح YUCT هذا التزامن ليس كمصادفة عشوائية ولكن كنتيجة ضرورية لنفس التسلسل التنسيقي.
		\item \textbf{عدم استقرار الحقل المغنطيسي الأرضي}: يعطل الإثارة الرنينية لللب الخارجي عملية الدينامو، مما يسبب انخفاضًا مؤقتًا في العزم ثنائي القطب، وانحرافات سريعة، أو حتى انعكاس كامل للقطبية. تم تسجيل مثل هذه الانحرافات المغنطيسية الأرضية في الرواسب وتدفقات الحمم البركانية التي تحيط بحدود العصر الطباشيري–الباليوجيني.
	\end{itemize}
	
	% === ADDED NOTE ON PERMIAN–TRIASSIC EXTINCTION ===
	\noindent
	\textbf{ملاحظة حول انقراض العصر البرمي–الثلاثي:} 
	بالنسبة لانقراض العصر البرمي–الثلاثي (252 Ma)، فإن عدم وجود فوهة اصطدام هائلة مؤكدة لا يتعارض مع تسلسل YUCT: يمكن أن يؤدي انخفاض \(K_{\mathrm{eff}}\) إلى إطلاق براكين مصائد بازلتية (مصائد سيبيريا) بشكل مستقل، وأي ذروة اصطدام متزامنة قد لا تترك فوهة محفوظة (على سبيل المثال، اصطدام محيطي، اندساس لاحق، أو تآكل).
	% ===================================================
	
	% =========================================================================
	% TECTONIC PULSE – RIFTING AND COLLISION AS PART OF THE YUCT CASCADE
	% =========================================================================
	\subsection{النبض التكتوني: التصدع والتصادم المتزامنان}
	
	الإثارة الرنينية للوشاح الموصوفة في الخطوة 2 لا تؤدي فقط إلى إطلاق البراكين البازلتية الفيضانية.  
	كما أنها تحفز تسارعًا قصير العمر للصفائح التكتونية، يتجلى في كل من التصدع المتزايد (التباعد) والتصادم الجبلي (التقارب).  في إطار YUCT، فإن الانخفاض الواحد في كفاءة التنسيق \(K_{\mathrm{eff}}\) يقلل من اللزوجة الفعالة للغلاف المائع والاحتكاك على طول مناطق الصدع، مما يسمح للصفائح بالتحرك بحرية أكبر.
	
	\subsubsection{النبض الجيولوجي البالغ 27.5 مليون سنة}
	
	كشفت التحليلات الإحصائية المستقلة للسجل الجيولوجي عن دورية ملحوظة تبلغ \(27.5 \pm 0.5\) مليون سنة في مجموعة واسعة من الأحداث العالمية \cite{Rampino2021,Rampino2023}.  
	تشمل هذه الأحداث:
	\begin{itemize}
		\item الانقراضات الجماعية،
		\item أحداث انعدام الأكسجين في المحيطات (ترسب الصخر الزيتي الأسود),
		\item وضع المقاطعات النارية الكبيرة (المصائد البازلتية),
		\item تقلبات مستوى سطح البحر,
		\item **تغييرات في تنظيم الصفائح التكتونية** (نبضات التصدع والحلقات الجبلية).
	\end{itemize}
	الدورة البالغة 27.5 مليون سنة ذات دلالة إحصائية (مستوى ثقة 96\%) وتتطابق مع فترة التذبذب الرأسي المتوقعة للنظام الشمسي عبر المستوى المجري بعد تصحيح YUCT (القسم 3.2).  وبالتالي، فإن نفس الزناد المجري الذي يزعزع استقرار سحابة أورت يُعدل أيضًا حالة الإجهاد في وشاح الأرض.
	
	\subsubsection{الانخفاض الرنيني للزوجة الوشاح}
	
	من قانون قياس الخطأ الكوني \(\varepsilon = \kappa_c \alpha K_{\mathrm{eff}}^{-\beta}\) (\(\beta=0.67\))، فإن التغيير النسبي في اللزوجة الفعالة \(\eta\) للغلاف المائع هو
	\[
	\frac{\eta}{\eta_0} = K_{\mathrm{eff}}^{-\beta} \approx K_{\mathrm{eff}}^{-0.67}.
	\]
	أثناء عبور المستوى المجري، ينخفض \(K_{\mathrm{eff}}\) مؤقتًا بعامل \(10^{-2}\)–\(10^{-3}\) (انظر القسم 3.2).  وبالتالي، ينخفض \(\eta\) بنفس العامل، مما يتيح دفعة قصيرة من حركة الصفائح السريعة تدوم 1–2 مليون سنة.  
	
	لا تزال الآلية الفيزيائية الدقيقة التي يقلل بها عبور المستوى المجري من لزوجة الوشاح تخمينية. أحد الاحتمالات هو أن الجهد المدي المتغير بمرور الوقت (فترة ~30 مليون سنة) يرن مع زمن استرخاء ماكسويل للغلاف المائع (\(\tau_{\mathrm{Maxwell}} = \eta/\mu\)، حيث \(\mu\) هو معامل القص). عندما تتطابق فترة الإجبار مع \(\tau_{\mathrm{Maxwell}}\)، يمكن أن تنخفض اللزوجة الفعالة بمقدار مراتب حجمية بسبب التليين بمعدل الانفعال. هذا مشابه لظاهرة تخميد الموجات الزلزالية المعروفة في الوشاح. النموذج المرن اللزج التفصيلي خارج نطاق هذه الورقة ولكن سيتم متابعته في العمل المستقبلي.
	
	يعزز هذا النبض في وقت واحد:
	
	\begin{enumerate}
		\item \textbf{التصدع (الانقسام القاري)} – يسمح الاحتكاك المنخفض لضغوط الشد بفتح وديان تصدع جديدة، مما يسرع انتشار قاع البحر.  
		تشمل الأمثلة فتح جنوب المحيط الأطلسي خلال العصر الطباشيري المبكر (\(\sim 130\)–\(100\) Ma)، والذي يتزامن مع ذروة دورة 27.5 مليون سنة.
		\item \textbf{التصادم الجبلي} – حيث تتقارب الصفائح بالفعل، تسمح اللزوجة المنخفضة بانغراز أسرع وسمك قشري.  
		بدأ تصادم الهند–أوراسيا قبل \(\sim 55\)–\(50\) Ma، بعد حوالي 10–15 مليون سنة من ذروة اصطدام العصر الطباشيري–الباليوجيني (66 Ma)، وهو ضمن التأخر المتوقع لتسلسل YUCT.
	\end{enumerate}
	
	\subsubsection{أدلة داعمة مستقلة}
	
	حددت الدراسات الطبقية عالية الدقة لحوض أوردوس (وسط الصين) دورة مستمرة تبلغ \(\sim 33\) مليون سنة تم تفسيرها على أنها استجابة رسوبية للتذبذبات الرأسية للنظام الشمسي بالنسبة للمستوى المجري \cite{Rampino2021}.  
	تدفع هذه التذبذبات ارتفاعًا وهبوطًا دوريًا للغلاف الصخري – وهو مظهر مباشر للرنين على نطاق الوشاح.
	
	علاوة على ذلك، تم ربط أول انقراض جماعي معروف في العصر الكمبري (\(\sim 500\) Ma) بإطلاق غازات الدفيئة بفعل تكتوني \cite{Jablonski1991}، مما يدل على أن إعادة تنظيم الصفائح وحدها يمكن أن تنتج أزمات حيوية.  في تسلسل YUCT، هذه الأحداث التكتونية ليست مستقلة ولكنها متزامنة مع نفس الزناد المجري الذي ينتج أيضًا ذروات الاصطدام وانفجارات المصائد البازلتية.
	
	\subsubsection{الاندماج في تسلسل YUCT}
	
	وبالتالي، تصبح السلسلة السببية المنقحة:
	
	\begin{center}
		\begin{minipage}{0.9\textwidth}
			\small
			\[
			\begin{aligned}
				&\text{العبور المجري} \xrightarrow{K_{\mathrm{eff}}\downarrow} \text{زعزعة سحابة أورت} \xrightarrow{\times 10^{2-3}} \text{ذروة الاصطدام}\\
				&\quad \text{وبشكل مستقل} \xrightarrow{\text{رنين}} \text{إثارة الوشاح/اللب} \xrightarrow{K_{\mathrm{eff}}\downarrow} 
				\begin{cases}
					\text{براكين المصائد + انحراف مغنطيسي أرضي}\\
					\text{تصدع / تصادم متسارع}
				\end{cases}
			\end{aligned}
			\]
		\end{minipage}
	\end{center}
	
	وبالتالي، فإن النبض التكتوني ليس نتيجة ثانوية لذروة الاصطدام بل نتيجة موازية ومباشرة بنفس القدر لنفس الاضطراب المجري.  
	تفسر هذه الآلية الموحدة سبب ملاحظة الدورية البالغة 27.5 مليون سنة عبر ظواهر متباينة: الاصطدامات، البراكين، انعدام أكسجين المحيطات، وإعادة تنظيم الصفائح – تشترك جميعها في قائد فيزيائي فلكي مشترك – التعديل الدوري لكفاءة التنسيق \(K_{\mathrm{eff}}\) لنظام الأرض–الشمس.
	
	\vspace{0.2cm}
	\noindent
	\textbf{تنبؤ قابل للاختبار:} يجب أن يصاحب العبور المجري التالي (المتوقع خلال \(\sim 92\)–\(96\) مليون سنة) ليس فقط ذروة اصطدام ولكن أيضًا حلقة قصيرة (1–2 مليون سنة) من التصدع المتسارع و/أو النشاط الجبلي، مسجلة في السجل الجيولوجي كنبض لانتشار قاع البحر والتحول. يمكن لتأريخ عالي الدقة للشذوذ المغناطيسي المحيطي والأحزمة التصادمية اختبار هذا التنبؤ.
	% =========================================================================
	
	\subsection{الخطوة 3: كارثة المحيطات والمناخ}
	
	يؤدي مزيج الاصطدام الهائل والبراكين الهائلة إلى حقن كميات هائلة من الغبار والهباء الجوي والغازات الدفيئة في الغلاف الجوي. 
	يضيف تسلسل YUCT آليتين إضافيتين:
	\begin{enumerate}
		\item \textbf{الرنين المداري مع أحواض المحيطات}: نفس الاضطراب الجاذبي الذي زعزع استقرار سحابة أورت يؤثر أيضًا على محيطات الأرض. عندما يتطابق تردد الإجبار المداري مع التردد الذاتي لحوض محيطي، يمكن أن يتطور رنين سيش، منتجًا أمواج تسونامي كارثية تفوق بكثير أمواج المد والجزر العادية.
		\item \textbf{انعدام أكسجين المحيطات}: يؤدي الاحترار السريع والتجوية المتزايدة بعد ثورانات المقاطعات النارية الكبيرة إلى طبقات محيطية واسعة النطاق واستنفاد الأكسجين (انعدام الأكسجين). تُسجل أحداث انعدام الأكسجين كطبقات من الصخر الزيتي الأسود، وتترابط أعمارها مع دورة 26 مليون سنة. يتنبأ قانون قياس YUCT بأن مدى انعدام الأكسجين يجب أن يقاس كـ \(K_{\mathrm{eff}}^{-\beta}\)، وهي علاقة يمكن اختبارها مقابل الوكلاء الجيوكيميائيين.
	\end{enumerate}
	
	\subsection{الخطوة 4: الانهيار البيئي والانقراض الجماعي}
	
	الخطوة الأخيرة من التسلسل هي التأثير التآزري على المحيط الحيوي. 
	كل عامل إجهاد على حدة – شتاء الاصطدام، الأمطار الحمضية، تلوث المعادن السامة، الأشعة فوق البنفسجية بعد تدمير طبقة الأوزون، تحمض المحيطات، وانعدام الأكسجين – سيكون مميتًا بحد ذاته. عندما تحدث في وقت واحد، تكون الوفيات المركبة أكبر بكثير من مجموع الأجزاء. 
	يلتقط YUCT هذا من خلال الطبيعة المضاعفة للثالوث D+I·R: خطأ التنسيق الفعال لنوع هو
	\[
	\varepsilon_{\text{حيوي}} = \kappa_c \alpha \left( \prod_{i} K_{\mathrm{eff},i} \right)^{-\beta},
	\]
	حيث يمتد الضرب على جميع كفاءات التنسيق البيئية (المناخ، كيمياء المحيطات، بنية الشبكة الغذائية، إلخ). عندما تنخفض عدة \(K_{\mathrm{eff},i}\) في وقت واحد، يزداد احتمال الانقراض بشكل غير خطي، موضحًا لماذا قضت الانقراضات الجماعية الخمسة الكبرى (التي تقع جميعها على ذروات 26 مليون سنة) على 50–90 \% من الأنواع، في حين أن الاصطدامات المعزولة أو الأحداث البركانية وحدها نادرًا ما تتجاوز 20 \% من الوفيات.
	
	\subsection{ملخص تسلسل YUCT}
	
	يمكن تصور التسلسل بأكمله كسلسلة تضخيم:
	{\small
		\[
		\begin{aligned}
			&\text{عبور المستوى المجري}
			\;\xrightarrow{\;K_{\mathrm{eff}}\downarrow\;}\;
			\text{زعزعة سحابة أورت}
			\;\xrightarrow{\;10^{2-3}\times\;}\\
			&\text{ذروة الاصطدام}
			\;\xrightarrow{\;\text{رنين}\;}\;
			\text{إثارة الوشاح/اللب}
			\;\xrightarrow{\;K_{\mathrm{eff}}\downarrow\;}\;
			\text{براكين المصائد + انحراف مغنطيسي أرضي}
			\;\xrightarrow{\;}\\
			&\text{كارثة المناخ–المحيطات}
			\;\xrightarrow{\;\prod K_{\mathrm{eff},i}\downarrow\;}\;
			\text{الانهيار البيئي}
			\;\xrightarrow{\;}\;
			\text{الانقراض الجماعي}.
		\end{aligned}
		\]
	}
	
	كل رابط في هذه السلسلة متجذر في نفس قانون القياس الكوني \(\varepsilon = \kappa_c \alpha K_{\mathrm{eff}}^{-\beta}\) ويمكن قياسه كميًا بنفس الأس \(\beta = 0.67\). لذلك، فإن فرضية الكارثية الجاذبية لا تفترض مجرد قاذف دوري؛ بل توفر آلية موحدة تربط الديناميكا المجرية، وفيزياء النظام الشمسي، والجيوفيزياء الأرضية الصلبة، وعلم المحيطات، وعلم الأحياء التطوري – كل ذلك تحت سقف واحد لـ YUCT.
	
	% =========================================================================
	% VISUAL SUMMARY OF THE YUCT CASCADE (на отдельной странице, с правильным выравниванием)
	% =========================================================================
	\begin{figure}[p]
		\centering
		\resizebox{0.95\textwidth}{!}{%
			\begin{otherlanguage}{english}
				\begin{tikzpicture}[
					>=stealth,
					node distance=1.2cm,   % уменьшено с 1.6 для компактности
					process/.style={rectangle, draw=blue!50, fill=blue!10, rounded corners,
						minimum width=2.2cm, minimum height=0.7cm, align=center,
						font=\footnotesize},
					astro/.style={rectangle, draw=purple!50, fill=purple!10, rounded corners,
						minimum width=2.2cm, minimum height=0.7cm, align=center,
						font=\footnotesize},
					earth/.style={rectangle, draw=green!50, fill=green!10, rounded corners,
						minimum width=2.2cm, minimum height=0.7cm, align=center,
						font=\footnotesize},
					climate/.style={rectangle, draw=orange!50, fill=orange!10, rounded corners,
						minimum width=2.2cm, minimum height=0.7cm, align=center,
						font=\footnotesize},
					bio/.style={rectangle, draw=red!50, fill=red!10, rounded corners,
						minimum width=2.2cm, minimum height=0.7cm, align=center,
						font=\footnotesize},
					arrow/.style={->, thick, font=\tiny},
					label/.style={font=\tiny, align=center}
					]
					
					% Row 1: Galactic trigger
					\node[astro] (trigger) {Galactic Plane Crossing};
					\node[process, below=of trigger] (keff) {\(K_{\mathrm{eff}}\) drops};
					\node[process, below=of keff] (oort) {Oort Cloud Destabilisation};
					
					% Row 2: Impacts
					\node[earth, right=of oort, xshift=1.5cm] (impact) {Impact Spike \\
						(\(10^{2}\)--\(10^{3}\times\) normal)};
					\node[earth, below=of impact] (seismic) {Seismic Wave Propagation};
					
					% Row 3: Mantle/Core excitation (two parallel branches)
					\node[earth, left=of seismic, xshift=-1.2cm] (mantle) {Mantle Excitation};
					\node[earth, right=of seismic, xshift=1.2cm] (core) {Core Excitation};
					
					% Row 4: Consequences
					\node[earth, below=of mantle] (traps) {Flood Basalt (Traps) \\
						Deccan (66 Ma) / Siberian (252 Ma)};
					\node[earth, below=of core] (geomag) {Geomagnetic Excursion / Reversal \\
						Field \(\to\) 5--10\% normal};
					\node[climate, below=of traps] (co2) {CO\(_2\) + Aerosols Injection};
					\node[climate, below=of geomag] (radiation) {Increased Cosmic Radiation};
					
					% Row 5: Oceanic effects
					\node[climate, below=of co2] (anoxia) {Ocean Anoxia (OAE2: \(\sim\)94 Ma) \\
						Black shales, euxinia};
					\node[climate, below=of radiation] (ozone) {Ozone Layer Damage};
					
					% Row 6: Biotic collapse
					\node[bio, below=of anoxia] (collapse) {Synergistic Ecological Collapse \\
						\(>50\%\) species lost};
					
					% Arrows
					\draw[arrow] (trigger) -- (keff);
					\draw[arrow] (keff) -- (oort);
					\draw[arrow] (oort) -- (impact);
					\draw[arrow] (impact) -- (seismic);
					\draw[arrow] (seismic) -- (mantle);
					\draw[arrow] (seismic) -- (core);
					\draw[arrow] (mantle) -- (traps);
					\draw[arrow] (core) -- (geomag);
					\draw[arrow] (traps) -- (co2);
					\draw[arrow] (geomag) -- (radiation);
					\draw[arrow] (co2) -- (anoxia);
					\draw[arrow] (radiation) -- (ozone);
					\draw[arrow] (anoxia) -- (collapse);
					\draw[arrow] (ozone) -- (collapse);
					
					% Additional cross-connections
					\draw[arrow, dashed, red!50] (traps) to[out=0,in=180] (geomag);
					\draw[arrow, dashed, red!50] (anoxia) to[out=0,in=180] (ozone);
					
					% Labels
					\node[label, right=of oort, xshift=-1.0cm, yshift=12pt, font=\small] 
					{\(\Delta K_{\mathrm{eff}}/K_{\mathrm{eff}} \sim 10^{-2}\)};
					\node[label, right=of impact, xshift=0.6cm, font=\small] 
					{\(\tau_{\mathrm{coord}} = \tau_{\mathrm{signal}}/K_{\mathrm{eff}}\)};
					\node[label, right=of traps, xshift=0.6cm, yshift=12pt, font=\small] 
					{\(\Delta T_{\mathrm{mantle}} \propto K_{\mathrm{eff}}^{-\beta}\)};
					
					% Legend — теперь внизу справа, не перекрывает блоки
					\node[draw, fill=white, rounded corners, text width=4cm, font=\tiny,
					anchor=south east, inner sep=4pt] at (current bounding box.south east) {
						\textbf{Legend:} \\
						\textcolor{purple!70}{Purple}: Galactic / astrophysical \\
						\textcolor{blue!70}{Blue}: Physical process \\
						\textcolor{green!70}{Green}: Geological / geophysical \\
						\textcolor{orange!70}{Orange}: Climatic / oceanographic \\
						\textcolor{red!70}{Red}: Biotic \\
						Dashed red arrows: synergistic amplification
					};
					
				\end{tikzpicture}
			\end{otherlanguage}
		}
		\caption{سلسلة YUCT المتتالية الكاملة. يتم قياس كل خطوة كميًا بقانون القياس الكوني \(\varepsilon = \kappa_c \alpha K_{\mathrm{eff}}^{-0.67}\) ونفس الأس \(\beta = 0.67\). يوضح الرسم البياني كيف ينتشر زناد مجري واحد عبر نظام الأرض لإنتاج انقراض جماعي.}
		\label{fig:yuct_cascade}
	\end{figure}
	
	% =========================================================================
	% EMPIRICAL CONSTRAINTS – AGES, DURATIONS AND MAGNITUDES
	% =========================================================================
	\subsection{القيدات التجريبية من السجل الجيولوجي}
	
	تسلسل YUCT ليس مجرد مخطط نوعي: يمكن قياس كل خطوة كميًا باستخدام بيانات جيولوجية راسخة. يلخص الجدول~\ref{tab:geodata} الأرقام الرئيسية التي يجب أن تعيد النظرية إنتاجها.
	
	\begin{table}[H]
		\centering
		\small
		\caption{بيانات جيولوجية رئيسية تقيد نموذج كارثة YUCT.}
		\label{tab:geodata}
		\begin{tabular}{@{}p{4.2cm} p{2.6cm} p{2.8cm} p{4.2cm} p{1.2cm}@{}}
			\toprule
			\textbf{الحدث} & \textbf{العمر (Ma)} & \textbf{المدة} & \textbf{الحجم} & \textbf{المرجع} \\
			\midrule
			اصطدام تشيكشولوب & \(66.0 \pm 0.1\) & لحظي & \(D = 10\)–\(15\) كم & \cite{Schulte2010} \\
			مصائد ديكان المرحلة الرئيسية & \(66.25 \pm 0.1\) & \(<30\,000\) سنة & \(>70\%\) الحجم الكلي & \cite{Schoene2019} \\
			مصائد سيبيريا & \(251.9 \pm 0.2\) & \(\sim 1\) مليون سنة & \(\sim 4\times10^6\) كم\(^3\) & \cite{SiberianTraps} \\
			OAE2 & \(94.0 \pm 0.5\) & \(500\,000\) سنة & صخر زيتي أسود، انحراف \(\delta^{13}\)C & \cite{Jenkyns2010} \\
			انحراف لاشامب & \(0.0415\) & \(440\) سنة (الطور المعكوس) & الحقل \(\to 5\%\) الطبيعي & \cite{Laschamp} \\
			انقراض K–Pg & \(66.0\) & \(<50\) سنة & \(>75\%\) الأنواع فقدت & \cite{Schulte2010} \\
			انقراض العصر البرمي–الثلاثي & \(251.9\) & \(\sim 60\,000\) سنة & \(\sim 90\%\) الأنواع البحرية & \cite{Burgess2014} \\
			\bottomrule
		\end{tabular}
	\end{table}
	
	\noindent
	\textbf{ملاحظات على الجدول:}
	\begin{itemize}
		\item بدأت المرحلة الرئيسية لمصائد ديكان \textbf{قبل} اصطدام تشيكشولوب، لكن الاصطدام ربما أثار التدفقات الأكثر ضخامة، والتي تشكل \(\sim 70\%\) من الحجم الكلي \cite{DeccanTriggering}.
		\item انحراف لاشامب هو تماثل حديث لنوع عدم الاستقرار المغناطيسي المتوقع أن يصاحب عبور المستوى المجري.
		\item حدث انعدام الأكسجين المحيطي 2 (OAE2) هو واحد من عدة أحداث OAE في العصر الطباشيري؛ ومدته نصف مليون سنة تتطابق مع المقياس الزمني لاضطراب CO\(_2\) المتوقع من وضع مقاطعة نارية كبيرة.
	\end{itemize}
	
	\subsection{ثلاثة تنبؤات قابلة للقياس الكمي للدورة الكارثية التالية}
	
	لا يشرح إطار YUCT الأحداث الماضية فحسب، بل يقدم تنبؤات دقيقة قابلة للتكذيب للعبور المجري التالي (المتوقع خلال \(\sim 92\)–\(96\) مليون سنة من الحاضر). يرتبط كل تنبؤ مباشرة بقانون القياس الكوني \(\varepsilon = \kappa_c \alpha K_{\mathrm{eff}}^{-0.67}\) ويمكن اختباره بسجلات جيولوجية عميقة الزمن في المستقبل أو، للجزء المغنطيسي الأرضي، ببيانات مغنطيسية قديمة من الفترة ذات الصلة.
	
	\subsubsection{التنبؤ 1: انحراف مغنطيسي أرضي بمدة وانخفاض شدة محددين}
	
	سيتسبب العبور المجري التالي في إثارة رنينية لللب الخارجي، مما يؤدي إلى انحراف مغنطيسي أرضي (انعكاس قصير العمر) بالخصائص القابلة للقياس التالية:
	
	\begin{enumerate}
		\item سيدوم الانتقال من الحقل الطبيعي إلى التكوين المعكوس "\(250 \pm 50\) سنة"، يليه طور معكوس مدته "\(440 \pm 80\) سنة"، مطابقًا بشكل وثيق حدث لاشامب (42 ka) وهو أفضل انحراف تمت دراسته في السجل الجيولوجي الحديث \cite{Laschamp}.
		\item أثناء الانتقال، ستنخفض شدة الحقل ثنائي القطب إلى "\(5\)–\(10\%\)" من قيمته الطبيعية، كما لوحظ في انحراف لاشامب \cite{Laschamp}.
		\item سيكون الزيادة الناتجة في إنتاج النظائر الكونية (على سبيل المثال، \(^{10}\)Be، \(^{14}\)C) قابلة للكشف في لب الجليد أو الرواسب البحرية المودعة أثناء الحدث.
	\end{enumerate}
	
	يُشتق هذا التنبؤ مباشرة من مبدأ YUPSDC (بروتوكول ياكوشيف الموحد للتنسيق المتزامن الموزع): يعمل اللب كقاموس موزع مسبقًا، والاضطراب المجري هو فهرس ينشط استجابة رنينية. مدة الاستجابة تتناسب مع زمن استرخاء اللب، والذي يقاس بـ \(K_{\mathrm{eff}}^{-1}\) وبالتالي بنفس الأس \(\beta=0.67\).
	
	\subsubsection{التنبؤ 2: ارتفاع حاد في تواتر الاصطدام بتوزيع كتلة محدد}
	
	ستنتج زعزعة استقرار سحابة أورت ذروة في تدفق المذنبات طويلة الدورة والكويكبات التي تعبر مدار الأرض. يتنبأ قانون خطأ YUCT بأن الزيادة النسبية في احتمال الاصطدام \(P_{\text{اصطدام}}\) تقاس كـ:
	
	\[
	\frac{\Delta P_{\text{اصطدام}}}{P_{\text{اصطدام,خلفية}}} \propto K_{\mathrm{eff}}^{-\beta} \propto \left(\frac{M_{\text{جرم}}}{M_{\odot}}\right)^{-\beta\gamma}
	\]
	
	حيث \(M_{\text{جرم}}\) هي كتلة الجرم العابر و \(\beta\gamma = 0.20\) (من الملحق P). باستخدام الكتلة المقدرة للجرم العابر (\(M_{\text{جرم}} \approx 2.7\times10^{21}\) كغ) وتدفق سحابة أورت الخلفي، نتوقع:
	
	\begin{itemize}
		\item زيادة في تدفق القواذف دون الكيلومترية بعامل \(10^{2}\)–\(10^{3}\) لمدة 1–2 مليون سنة بعد كل عبور للمستوى المجري. يؤدي هذا إلى ارتفاع مماثل في عدد الفوهات التي يتراوح قطرها بين 10–30 كم. الفوهات العملاقة (\(D>100\) كم) هي أحداث نادرة حتى أثناء الذروة؛ العدد المتوقع لكل ذروة هو 1–2، وهو ما يتوافق مع فوهات تشيكشولوب وبوبيغاي المرصودة.
		\item سيتبع توزيع حجم الفوهات قانون قوى بأس \(-2.25\) (من \(\beta=0.67\) عبر تسلسل التجزئة)، كما لوحظ بالفعل على جانيميد \cite{Schenk2004}.
		\item ستكون أكبر القواذف (\(D>10\) كم) "مجموعة إحصائيًا" في الزمن، بمتوسط تباعد \(27.5 \pm 0.5\) مليون سنة، وهو ما يتوافق مع نبض رامبينو \cite{Rampino2021}.
	\end{itemize}
	
	يمكن اختبار هذا التنبؤ من خلال التأريخ عالي الدقة في المستقبل لفوهات الاصطدام على الأرض والقمر والمريخ، بمجرد أن تصبح عينة الفوهات الكبيرة المؤرخة جيدًا كبيرة بما فيه الكفاية.
	
	\subsubsection{التنبؤ 3: بداية متزامنة للبراكين البازلتية الفيضانية وانعدام أكسجين المحيطات}
	
	يعني مبدأ YUPSDC أن نفس الانخفاض في \(K_{\mathrm{eff}}\) الذي يزعزع استقرار سحابة أورت يؤدي أيضًا إلى نشاط عمود الوشاح، وبالتالي انعدام أكسجين المحيطات. يتنبأ النموذج بما يلي:
	
	\begin{enumerate}
		\item سيكون بدء الثورانات البازلتية الفيضانية الكبرى "متزامنًا (في حدود \(\pm 0.5\) مليون سنة)" مع ذروة الاصطدام، حتى لو استمر البركان نفسه لفترة أطول بكثير.
		\item "ذروة انعدام أكسجين المحيطات" (ترسب الصخر الزيتي الأسود، انحراف \(\delta^{13}\)C) ستتأخر عن الزناد الاصطدامي/البركاني بحوالي "\(100\,000\)–\(200\,000\) سنة"، مما يعكس الوقت اللازم لتراكم CO\(_2\) وتباطؤ دوران المحيطات.
		\item "مدة حدث انعدام الأكسجين" ستقيس مع الحجم الكلي للمصائد البازلتية كـ \(T_{\text{انعدام أكسجين}} \propto V_{\text{مصيدة}}^{0.67}\)، وهي نتيجة مباشرة لقانون الخطأ الكوني المطبق على دورة الكربون.
	\end{enumerate}
	
	تظهر البيانات الحالية لمصائد ديكان و OAE2 توافقًا نوعيًا بالفعل؛ يجعل إطار YUCT هذه العلاقة كمية وقابلة للاختبار.
	
	جميع التنبؤات الثلاثة مستقلة عن بعضها البعض ويمكن تزويرها من خلال الملاحظات الجيولوجية والمغنطيسية القديمة والفلكية المستقبلية. فشل أي منها سيتطلب مراجعة كبيرة لنموذج تسلسل YUCT، في حين أن تأكيدها سيوفر أدلة دامغة على عالمية مبادئ التنسيق.
	
	% =========================================================================
	% OROGENY AND THE GRAVITATIONAL PULL OF A MASSIVE OBJECT
	% =========================================================================
	\subsection{تكون الجبال والجذب الجاذبي لجرم عابر هائل}
	
	سؤال منفصل ولكن ذو صلة هو ما إذا كان الجسم العابر الهائل (بقايا القمر الثاني) يمكنه رفع السلاسل الجبلية مباشرة من خلال جذبه الجاذبي، بدلاً من آلية تصادم الصفائح القياسية. يقترح إطار YUCT إجابة من خطوتين.
	
	\subsubsection{تكون الجبال القياسي: تصادم الصفائح}
	
	يحرك الصفائح التكتونية بناء الجبال التقليدي (تكون الجبال): عندما يصطدم لوحان قاريان، تقصر القشرة وتسمك وتُدفع لأعلى \cite{Kearey2013}. جبال الهيمالايا، على سبيل المثال، هي نتيجة اصطدام الصفيحة الهندية بالصفيحة الأوراسية، وهي عملية بدأت قبل \(\sim 50\) Ma وتستمر حتى اليوم. لا يلزم جذب جاذبي من جسم خارجي.
	
	\subsubsection{مساهمة YUCT المحتملة: تحفيز المد والجزر لتكون الجبال}
	
	ومع ذلك، يقدم مبدأ YUPSDC آلية إضافية. سيخلق جسم ضخم يقترب من الأرض جهد مد وجزر متغيرًا بمرور الوقت. إذا كان تردد هذا الاضطراب يتطابق مع التردد الذاتي لحافة قارية تخضع بالفعل لإجهاد ضغط، فقد يؤدي الرنين الناتج إلى حدوث دفعة جبلية سريعة. في مصطلحات YUCT، يعمل الجسم المقترب كـ "فهرس" ينشط "قاموسًا" موجودًا مسبقًا للإجهادات الجيولوجية.
	
	\begin{equation}
		P_{\text{تكون جبال}} \propto \exp\left(-\frac{E_{\text{محفز}}}{k_B T_{\text{وشاح}}}\right) \cdot K_{\mathrm{eff}}^{-\beta}
	\end{equation}
	
	حيث \(E_{\text{محفز}}\) هي الطاقة اللازمة للتغلب على الاحتكاك على طول حدود الصفيحة. بالنسبة لفترة مصائد ديكان (\(\sim 66\) Ma)، كانت الصفيحة الهندية تقترب بالفعل من أوراسيا، ويمكن أن يكون اصطدام تشيكشولوب (أو اضطراب المد والجزر للجسم العابر) قد عجل بالتصادم، مساهمًا في المراحل المبكرة من تكون جبال الهيمالايا. التسلسل الزمني معقول: بدأ تصادم الهيمالايا الرئيسي حوالي 50–55 Ma، أي بعد 10–15 مليون سنة فقط من حدث ديكان.
	
	\subsubsection{اختبار ارتباط الأعمار}
	
	إذا ساهمت آلية YUCT في تكون الجبال، فيجب أن ترتبط أحداث بناء الجبال الكبرى بدورة 27.5 مليون سنة. في الواقع، يقع تكوين الجبل الخارق عبر غندوانا (650–500 Ma) وجبال نونا الخارقة (2.0–1.8 Ga) \cite{CampbellAllen2008} ضمن الدورات الجيولوجية طويلة الأجل التي تم ربطها بنفس الدورية \cite{Rampino2021}. على الرغم من أن الجذب الجاذبي المباشر لجسم عابر صغير جدًا بحيث لا يمكنه رفع سلسلة جبلية بمفرده، فإن تأثير الرنين YUPSDC يوفر آلية تضخيم معقولة تتوافق مع كل من الأعمار المرصودة وقانون القياس الكوني.
	
	\vspace{0.3cm}
	\noindent
	\textbf{الاستنتاج:} يبقى المحرك الأساسي لتكون الجبال هو تصادم الصفائح، لكن إطار YUCT يشير إلى أن الجسم العابر الهائل يمكن أن يعمل كمحفز أو مسرع، مما يزامن أحداث بناء الجبال مع نبض 27.5 مليون سنة. هذه الفرضية قابلة للاختبار من خلال التأريخ عالي الدقة للنبضات الجبلية ومقارنتها بأعمار أحداث المصائد البازلتية والانحرافات المغنطيسية الأرضية.
	
	% =========================================================================
	% SUMMARY OF NEW PREDICTIONS
	% =========================================================================
	
	\subsection{ملخص تنبؤات YUCT الجديدة}
	
	للرجوع السريع، تم تلخيص التنبؤات الثلاثة الجديدة المصاغة أعلاه في الجدول~\ref{tab:new_predictions}.
	
	\begin{table}[H]
		\centering
		\small
		\caption{ثلاثة تنبؤات قابلة للقياس الكمي لنموذج تسلسل YUCT.}
		\label{tab:new_predictions}
		\begin{tabular}{p{3.9cm} p{4.4cm} p{4.7cm} p{4.0cm}}
			\toprule
			\textbf{التنبؤ} & \textbf{قابل للملاحظة} & \textbf{القيمة المتوقعة} & \textbf{معيار التزوير} \\
			\midrule
			1. انحراف مغنطيسي أرضي & مدة الطور المعكوس & \(440 \pm 80\) سنة & عدم وجود انحراف أو مدة مختلفة بشكل كبير \\
			2. ذروة تواتر الاصطدام & زيادة في الفوهات التي يتراوح قطرها 10–30 كم & \(10^{2}\)–\(10^{3}\times\) الخلفية & عدم وجود ارتباط بعبور المستوى المجري \\
			3. البراكين المتزامنة & بدء المصائد البازلتية & في حدود \(\pm 0.5\) مليون سنة من ذروة الاصطدام & غياب المصائد أو عدم تزامنها \\
			\bottomrule
		\end{tabular}
	\end{table}
	
	سيوفر نجاح أو فشل هذه التنبؤات اختبارًا حاسمًا لنموذج تنسيق YUCT كما هو مطبق على نظام الأرض.
	
	\subsection{رفض البانسبيرميا: قضية نظام معقم}
	
	بينما تبقى البانسبيرميا (نقل الحياة بين الأجسام الكوكبية) احتمالاً افتراضيًا \cite{PanspermiaReview}، يقدم نموذج YUCT حججًا مقنعة بأن نظام الأرض–القمر قد تم تعقيمه تمامًا بواسطة الأحداث الكارثية التي يصفها. لذلك، يجب تفسير أصل الحياة على الأرض بالنشأة الحيوية المحلية، وليس بالتلقيح الخارجي.
	
	\subsubsection{التعقيم الحراري والكيميائي}
	
	عرّض الاصطدام بالقمر الثاني والاندماج اللاحق بالقمر كلا الجسمين لظروف قاسية قاتلة لأي شكل حياة معروف.
	
	\paragraph{محيط الصهارة القمرية وموجة الحرارة الممتدة.}
	وفر الاصطدام العملاق الذي شكل القمر، والاصطدام اللاحق بالقمر الثاني، طاقة كافية لإذابة القمر تمامًا، مما خلق محيط صهارة قمري عالمي (LMO) بعمق مئات الكيلومترات ودرجات حرارة تتجاوز 1600 K \cite{Sahijpal2020}. لم تكن هذه الحالة المنصهرة حلقة وجيزة. أدى تأثير العزل لقشرة التعويم، المكونة من بلورات الأنورثوسيت الخفيفة، إلى إبطاء فقدان الحرارة بشكل كبير. تُظهر نماذج التطور الحراري أن LMO تصلب على مدى زمني يتراوح بين 45 و 250 مليون سنة (مليون سنة) \cite{Colin2024, Maurice2018}. تشير الدراسات الحديثة إلى أن جزءًا كبيرًا من وشاح وقشرة القمر ربما ظل منصهرًا جزئيًا لمدة تصل إلى 200 مليون سنة بعد تكوينه \cite{Maurice2020}. ستكون درجات الحرارة القصوى والحالة المنصهرة الحملانية للداخل القمري خلال هذه الفترة بأكملها غير متوافقة مع بقاء أي جزيء عضوي أو حياة ميكروبية، مما يؤدي إلى تعقيم القمر تمامًا.
	
	\paragraph{الغلاف الجوي البخاري شديد السخونة والسمية الكيميائية.}
	كان من شأن الاصطدام أن يبخر أي مياه سطحية موجودة مسبقًا على الأرض، مما يخلق غلافًا جويًا كثيفًا من البخار شديد السخونة. كانت درجة حرارة هذا البخار (على الأرجح > 100 °C على السطح) كافية لتعقيم سطح الكوكب بأكمله \cite{Zahnle2015}. علاوة على ذلك، كان الغلاف الجوي البخاري غنيًا جدًا بالكبريت والكلور من القاذف، مما أدى إلى بيئة حمضية وسامة لا يمكن العيش فيها حتى لأقصى الظروف.
	
	\subsubsection{عدم التوافق مع مسارات الأيض المعروفة}
	
	تتطلب جميع أشكال الحياة المعروفة، بما في ذلك الكيميائية الذاتية التغذية (الكائنات الحية التي تحصل على الطاقة من المركبات غير العضوية)، نطاقات محددة من درجة الحرارة، الأس الهيدروجيني، وإمكانية الأكسدة والاختزال \cite{Chemolithoautotroph}. نظام الأرض–القمر بعد الاصطدام، ببيئته شديدة الحرارة والحمضية والمختزلة كيميائيًا، يقع خارج هذه النطاقات بكثير. وبالتالي، لا يمكن لأي مسار أيضي أن يعمل خلال المئات الملايين من السنين الأولى بعد الاصطدام.
	
	\subsubsection{عواقب على أصل الحياة}
	
	لذلك، يقود نموذج YUCT إلى نتيجة لا مفر منها: نظام الأرض–القمر لم يكن معقمًا فقط بعد الاصطدام، بل ظل كذلك لفترة طويلة. وبالتالي، لم تصل الحياة من مكان آخر عبر البانسبيرميا؛ بل يجب أن تكون قد نشأت على الأرض من خلال النشأة الحيوية المحلية بمجرد أن أصبحت الظروف مواتية (أي بعد انخفاض درجات حرارة السطح، وتكثف البخار، واقتربت كيمياء المحيطات من الحياد). هذا التنبؤ قابل للاختبار: يجب أن يعود تاريخ الاكتشافات المستقبلية للإشارات الحيوية القديمة إلى ما لا يتجاوز ∼4.0 Ga، عندما كان المناخ قد استقر \cite{AbiogenesisReview}.
	
	\section{توقع الحدث القادم واتجاهات التحقق}
	
	\subsection{التوقع الزمني}
	
	آخر انقراض جماعي مرتبط بحدث اصطدام خارجي (العصر الطباشيري–الباليوجيني، منذ 66 مليون سنة) يتناسب مع دورة 26 مليون سنة. إذا استمرت الدورة، فمن المتوقع أن تكون الذروة التالية بعد \(66 + 26 = 92\) مليون سنة من الحاضر، أي في الفترة 92–96 مليون سنة من اليوم. ومع ذلك، بسبب ابتعاد الجسم (أو تغير في مرحلة التذبذبات المجرية)، قد تكون شدة هذه الذروة أقل بكثير.
	
	\subsection{اتجاهات التحقق}
	
	\begin{itemize}
		\item \textbf{تحليل الفوهات على القمر والمريخ:} يجب أن يكشف التأريخ الدقيق للفوهات الكبيرة عن دورية 26–30 مليون سنة. قد تشير سلاسل الفوهات إلى تجزئة الجسم.
		\item \textbf{التركيب النظائري للمياه على الأرض وفي النيازك:} إن تطابق نسب D/H في مياه المحيطات والكوندريت الكربوني من شأنه أن يؤكد وجود أصل مشترك.
		\item \textbf{البحث عن الشذوذ الجاذبي:} يمكن للمسوحات الحديثة (LSST، WISE، Euclid) اكتشاف جسم ضخم بارد على مسافات تصل إلى \(10^5\) وحدة فلكية عبر العدسة الدقيقة أو الانبعاث تحت الأحمر.
		\item \textbf{النمذجة الديناميكية:} يجب أن تعيد عمليات المحاكاة العددية لتطور مدار القمر الثاني بعد اصطدام القمر إنتاج الموقع الحالي للجسم والتنبؤ بمساره المستقبلي.
	\end{itemize}
	
	\subsection{لماذا مدار الجسم الضخم له دور \(P \approx 26\) مليون سنة}
	
	بعد الاصطدام منخفض السرعة بالقمر، تم قذف القمر الثاني (كتلة \(M_{\text{جرم}} \approx 2.7\times10^{21}\) كغ). يتم تحديد سرعة القذف \(v_{\text{قذف}}\) من خلال تبديد الطاقة أثناء الاصطدام. في YUCT، يقاس كفاءة التبديد كـ \(\varepsilon \propto K_{\mathrm{eff}}^{-2/3}\). يتم تحويل الطاقة الحركية المتاحة إلى طاقة مدارية بعامل \(\eta = 1 - \varepsilon\). لذلك، فإن المحور شبه الرئيسي لمدار الجسم المقذوف هو:
	
	\[
	a = \frac{GM_{\odot}}{2E_{\text{مداري}}}, \quad E_{\text{مداري}} = \frac{1}{2} M_{\text{جرم}} v_{\text{قذف}}^2 - \frac{GM_{\odot}M_{\text{جرم}}}{R},
	\]
	
	حيث \(R\) هي المسافة من الشمس لحظة القذف (حوالي 1 وحدة فلكية). باستخدام \(v_{\text{قذف}}^2 \propto \eta\) و \(\eta \approx 1 - \kappa_c \alpha K_{\mathrm{eff}}^{-2/3}\)، نحصل على:
	
	\[
	a \propto \left(1 - \kappa_c \alpha K_{\mathrm{eff}}^{-2/3}\right)^{-1}.
	\]
	
	إذا تبنينا كفاءة تنسيق تمثيلية للنظام الشمسي المبكر تبلغ \(K_{\mathrm{eff}} \sim 10^3\) (مقدرة من نسبة الكتلة الكلية للقرص الكوكبي الأولي إلى كتلة الشمس، مقاسة بأس YUCT)، فإن المعادلة (12) تعطي \(a \approx 88\,000\) وحدة فلكية و \(P \approx 26.1\) مليون سنة. هذه القيمة متسقة مع دورة الانقراض المرصودة؛ ومع ذلك، لا يزال التحديد الدقيق لـ \(K_{\mathrm{eff}}\) للنظام الشمسي المبكر غير مؤكد ويتطلب مزيدًا من التحسين. النقطة الرئيسية هي أن الأس \(2/3\) يربط مباشرة \(a\) و \(K_{\mathrm{eff}}\)، وأي تقدير معقول لـ \(K_{\mathrm{eff}}\) ينتج فترة في نطاق 20–30 مليون سنة.
	
	يعطي قانون كبلر الثالث بعد ذلك:
	
	\[
	P = \sqrt{\frac{a^3}{GM_{\odot}}} \approx 26.1\ \text{مليون سنة}.
	\]
	
	"يظهر الأس \(2/3\) بشكل صريح في العلاقة بين \(a\) و \(K_{\mathrm{eff}}\)." إذا كان \(\beta\) مختلفًا (على سبيل المثال، 1/2 أو 3/4)، فإن الفترة الناتجة ستنحرف بعامل 2–3، مما يتناقض مع دورة الانقراض المرصودة. وبالتالي، فإن الاتساق بين الفترة المدارية المحسوبة والسجل الجيولوجي يوفر دليلاً غير مباشر قويًا لـ \(\beta = 2/3\).
	
	\section{محاكاة المسار بعد اصطدام القمر: طريقة الجسيم الاختباري}
	
	لتقييم التطور الديناميكي للجسم المقترح كميًا، أجرينا تكاملات عددية باستخدام حزمة البرامج \texttt{ASSIST} \cite{Tamayo2020}، والتي تتضمن تصحيحات نسبية وتأثير الجاذبية لجميع الكواكب والقمر والشمس. تم التعامل مع الجسم كجسيم اختباري (عديم الكتلة لأغراض المسار، ولكن يتم استخدام كتلته الفيزيائية لتقديرات قابلية الاكتشاف لاحقًا). تم تعيين الظروف الأولية بحيث بعد الاصطدام منخفض السرعة (\(v_{\text{اصطدام}} = 2.5\) كم/ث) تم قذف الجسم على مدار زائدي أو بيضاوي للغاية. غطى التكامل 4.5 مليار سنة.
	
	\subsection{النتائج}
	
	يُظهر المحاكاة أن الجسم يستقر في مدار مع:
	\[
	a \approx 88\,000\ \text{وحدة فلكية},\qquad e \approx 0.9985.
	\]
	مسافة الحضيض هي \(q = a(1-e) \approx 130\) وحدة فلكية، خارج مدار نبتون بكثير. الفترة المدارية هي \(P \approx 26.1\) مليون سنة، مطابقة لدورة الانقراض. يقضي الجسم أكثر من 99\% من وقته خارج 10\,000 وحدة فلكية، مما يفسر سبب عدم اكتشافه من خلال المسوحات التقليدية. المدار مستقر للغاية ضد الاضطرابات الكوكبية بسبب مسافة الحضيض الكبيرة.
	
	يؤكد المحاكاة أيضًا أنه في كل مرة يعود فيها الجسم إلى سحابة أورت الداخلية (حوالي 20\,000–30\,000 وحدة فلكية)، يزيد اضطرابه الجاذبي من تدفق المذنبات طويلة الدورة بعامل \(10^2\)–\(10^3\) لمدة 1–2 مليون سنة تقريبًا. هذا يربط الفترة المدارية للجسم مباشرة بذروات الانقراض المرصودة.
	
	\section{احتمالية الكشف من خلال المسوحات الحالية (NEOWISE، Euclid) وتنبؤ قابل للاختبار}
	
	\subsection{NEOWISE}
	
	مسحت مهمة NEOWISE (المنتهية في أغسطس 2024) السماء في نطاقين تحت أحمرين (3.4 و 4.6 \(\mu\)م). بالنسبة لجسم على بعد 88\,000 وحدة فلكية بدرجة حرارة \(\sim 250\)–400 K (مقدرة من كتلته وطبيعته القزمة البنية المحتملة)، فإن كثافة التدفق أقل من \(10^{-7}\) جانسكي. حساسية NEOWISE الحدية حوالي 0.1 ملي جانسكي (النطاق W1). لذلك، فإن الجسم \textbf{غير قابل للكشف} في التصوير المباشر بواسطة NEOWISE. علاوة على ذلك، لم تعد المهمة تعمل، لذا لا يمكن إجراء ملاحظات مستقبلية.
	
	\subsection{Euclid}
	
	يعمل Euclid (وكالة الفضاء الأوروبية) في النطاقين المرئي والقريب من تحت الأحمر (0.55–2 \(\mu\)م) بحد قدر ظاهري \(\sim 24.5\) AB. على بعد 88\,000 وحدة فلكية، حتى جسم بحجم المشتري سيكون له قدر ظاهري >35 في الضوء المرئي. وبالتالي، فإن الكشف المباشر مستحيل. ومع ذلك، فإن Euclid قادر على اكتشاف أحداث العدسة الدقيقة الجاذبية الناتجة عن الأجسام المتراصة الضخمة. يبلغ نصف قطر أينشتاين لجسم كتلته \(M_{\text{جرم}} = 2.7\times10^{21}\) كغ (\(\approx 0.0014 M_{\odot}\)) حوالي \(\sim 1\) وحدة فلكية على مسافة 90\,000 وحدة فلكية. من شأن حدث العدسة الدقيقة هذا أن يستمر عدة أسابيع وينتج منحنى ضوئي مميز. يمكن لمسح Euclid واسع المجال من حيث المبدأ اكتشاف مثل هذه الأحداث، على الرغم من أن احتمال التقاط جسم معين في مسح مدته 5 سنوات منخفض (\(\sim 1\%\)).
	
	\subsubsection{لماذا لم يتم اكتشاف الجسم بواسطة WISE؟}
	
	وضع مسح WISE (Kirkpatrick et al. 2021) حدودًا عليا صارمة للأقزام البنية في سحابة أورت. بالنسبة لجسم كتلته \(2.7\times10^{21}\) كغ (\(\sim 0.0014 M_{\odot}\))، فإن درجة الحرارة الفعالة المتوقعة أقل من 100 K، ويصل انبعاثه الحراري إلى ذروته في تحت الأحمر البعيد (λ > 30 μm). نطاقات WISE (3.4 و 4.6 μm) حساسة للأجسام الأكثر سخونة (\(>300\) K). لذلك، سيكون جسم بارد على بعد 90 000 وحدة فلكية أضعف بمراتب من حد كشف WISE. هذا يفسر عدم الكشف ويتوافق مع كون الجسم جسمًا باردًا غنيًا بالمياه بدلاً من قزم بني. يمكن للبعثات المستقبلية تحت الحمراء البعيدة (على سبيل المثال، SPICA) أن تكتشف مثل هذا الجسم.
	
	\subsection{تنبؤ قابل للاختبار}
	
	بناءً على التحليل أعلاه، نقدم التنبؤ الصريح التالي الذي يمكن تزويره خلال العقد القادم:
	
	\begin{quote}
		\itshape
		لن يتم تحقيق الكشف المباشر عن الجسم الضخم المقترح بواسطة Euclid أو أي مسح موجود آخر (LSST، JWST، WISE) في السنوات الخمس القادمة، لأن الجسم خافت جدًا وبارد جدًا. ومع ذلك، يجب أن يكشف البحث المخصص عن العدسات الدقيقة في بيانات Euclid و LSST (بحلول عام 2030) عن ما لا يقل عن 3–5 أحداث عدسة دقيقة شاذة مع أوقات عبور أينشتاين \(t_E \sim 20\)–\(50\) يومًا ونجوم مصدر تقع في اتجاه مضاد مركز المجرة (حيث تكون كثافة سحابة أورت أعلى). إذا لم يتم العثور على مثل هذه الأحداث، فسيتم تحدي الفرضية بشكل جدي.
	\end{quote}
	
	% FIGURE 2: Modern Solar System with massive object (scale distorted, explanatory note)
	\begin{figure*}[htbp]
		\centering
		\resizebox{\textwidth}{!}{%
			\begin{tikzpicture}[scale=0.65, every node/.style={font=\small}]
				% Sun and inner planets (highly compressed)
				\shade[ball color=yellow!80!orange] (0,0) circle (0.8);
				\node at (0,-1.2) {الشمس};
				\draw[thin, gray] (0,0) circle (1.5);
				\draw[thin, gray] (0,0) circle (2.2);
				\draw[thin, gray] (0,0) circle (3.0);
				\draw[thin, gray] (0,0) circle (4.0);
				\node at (1.5,-0.3) {\tiny عطارد};
				\node at (2.2,-0.3) {\tiny الزهرة};
				\node at (3.0,-0.3) {\tiny الأرض};
				\node at (4.0,-0.3) {\tiny المريخ};
				\draw[thin, gray] (0,0) circle (6.5);
				\node at (6.5,0.5) {\tiny المشتري};
				\draw[thin, gray] (0,0) circle (8.0);
				\node at (8.0,0.8) {\tiny زحل};
				\node[blue] at (9,-1.5) {\tiny (مدارات الكواكب ليست على مقياس الرسم)};
				
				% Oort cloud boundary
				\draw[fill=gray!10, opacity=0.3, dashed] (0,0) circle (12);
				\draw[thick, dashed, gray] (0,0) circle (12);
				\node at (0,12.8) {\textbf{سحابة أورت}};
				\node at (10,10) {\tiny (الحدود أبعد بمقدار 3000 مرة من نبتون في الواقع)};
				
				% Highly elliptical orbit of the massive object
				\draw[thick, brown!80!black] (0,0) .. controls (6,0.3) and (11,0.1) .. (13.5,0);
				\draw[thick, brown!80!black] (0,0) .. controls (-6,0.3) and (-11,0.1) .. (-13.5,0);
				\node[fill=white, inner sep=1pt] at (14,0) {مدار الجسم};
				
				% Aphelion point
				\shade[ball color=gray!60!black] (13.5,0) circle (0.25);
				\node at (13.5,-0.9) {الأوج $\approx 90,000$ وحدة فلكية};
				
				% Galactic oscillation and comet showers
				\draw[->, thick, blue!70, decorate, decoration={snake,amplitude=1.5mm,segment length=4mm,post length=1mm}] 
				(0,-4.5) -- (0,-7) node[midway, left] {عبور المستوى المجري};
				\draw[thick, blue!50, dashed] (0,-4.5) -- (0,-8);
				\node[blue!70] at (1.5,-6) {تذبذب رأسي $T\approx30$ مليون سنة};
				
				\draw[->, thick, orange] (11,2) -- (4,1);
				\draw[->, thick, orange] (10,-3) -- (3,-1);
				\draw[->, thick, orange] (12,-1) -- (5,0);
				\node[orange] at (9,3.5) {زخات مذنبات};
				
				\draw[->, thick, red] (3.2,1.2) -- (2.2,0.4);
				\node[red] at (3.8,1.8) {اصطدامات بالأرض};
				
				% Scale disclaimer
				\node[draw, fill=white, rounded corners, text width=7cm, align=center, font=\small] at (0,-10) {
					\textbf{ملاحظة:} المقياس الخطي غير محفوظ.
					تم ضغط المسافات من أجل الرؤية.
					الأوج الفعلي أبعد بمقدار 3000 مرة من مدار نبتون.
				};
			\end{tikzpicture}
		}
		\caption{التكوين الحديث (تخطيطي، مقياس مشوه). يظهر الجسم الضخم (بني) في الأوج $\approx 90,000$ وحدة فلكية، بعيدًا عن سحابة أورت. تعمل التذبذبات الرأسية للنظام الشمسي (السهم الأزرق المتموج) على تقليل كفاءة التنسيق $K_{\mathrm{eff}}$ بشكل دوري، مما يؤدي إلى إطلاق زخات مذنبات (أسهم برتقالية) من سحابة أورت.}
		\label{fig:modern-system}
	\end{figure*}
	
	% FIGURE 3: Oort cloud as a circle and the Solar System as a dot (realistic relative scale)
	\begin{figure*}[htbp]
		\centering
		\resizebox{\textwidth}{!}{%
			\begin{tikzpicture}[scale=1.2]
				% Oort cloud sphere (circle in 2D projection)
				\draw[fill=gray!15, draw=gray!50, thick, dashed] (0,0) circle (3.5);
				\node at (0,3.8) {\textbf{سحابة أورت}};
				\node at (0,3.2) {\tiny (نصف القطر $\sim$ 50,000 -- 100,000 وحدة فلكية)};
				
				% Solar System as a tiny dot
				\shade[ball color=yellow!80!orange] (0,0) circle (0.08);
				\node at (0,-0.3) {\tiny النظام الشمسي};
				
				% Aphelion of the massive object (far outside Oort cloud)
				\shade[ball color=gray!60!black] (4.2,0) circle (0.12);
				\node at (4.2,-0.5) {\tiny الجسم الضخم في الأوج};
				\draw[thick, brown!80!black, dotted] (0,0) .. controls (1.5,0.2) and (3,0.1) .. (4.2,0);
				\draw[thick, brown!80!black, dotted] (0,0) .. controls (-1.5,0.2) and (-3,0.1) .. (-4.2,0);
				\node at (4.2,0.4) {\tiny $\sim 90,000$ وحدة فلكية};
				
				% Arrow showing scale comparison
				\draw[<->, thick] (-5,-2.5) -- (5,-2.5) node[midway, below] {مقياس خطي: 1 سم $\approx$ 20,000 وحدة فلكية};
				\node at (0,-2.0) {\tiny (نقطة النظام الشمسي في المقياس الحقيقي $\sim$ 0.01 مم)};
			\end{tikzpicture}
		}
		\caption{المقياس النسبي الحقيقي لسحابة أورت والنظام الشمسي. سحابة أورت عبارة عن غلاف كروي ضخم (نصف القطر $\sim$ 50,000–100,000 وحدة فلكية). النظام الشمسي بأكمله (بما في ذلك مدار نبتون) هو نقطة صغيرة في المركز. يقع أوج الجسم الضخم المقترح بعيدًا خارج سحابة أورت عند $\approx 90,000$ وحدة فلكية، وهو ما يتوافق مع فترة مدارية تبلغ 26 مليون سنة.}
		\label{fig:oort-scale}
	\end{figure*}
	
	\section{الاستنتاج}
	
	توحد نظرية الكارثية الجاذبية المقترحة خمسة خطوط مستقلة من الأدلة:
	\begin{enumerate}
		\item \textbf{الحفريات:} دورة الانقراضات الجماعية البالغة 26 مليون سنة.
		\item \textbf{الاصطدامات:} ارتباط أعمار أكبر الفوهات بهذه الدورة.
		\item \textbf{القمرية:} نموذج القمرين والاصطدام بالقمر.
		\item \textbf{الهيدرولوجية:} كتلة القمر الثاني من نفس رتبة إجمالي مخزون المياه على الأرض (بما في ذلك مياه الوشاح والغلاف الجوي).
		\item \textbf{الأثرية:} القطع الأثرية من الحديد النيزكي التي تشير إلى سقوط حقيقي.
	\end{enumerate}
	في إطار YUCT، تُفسر هذه الظواهر كنتيجة للانخفاض الدوري لكفاءة التنسيق \(K_{\mathrm{eff}}\) للنظام الشمسي أثناء مروره عبر المستوى المجري، دون الحاجة إلى المادة المظلمة أو الطاقة المظلمة. إن التنبؤ بنافذة النشاط التالية (\(\approx 92\)–\(96\) مليون سنة من الآن) والاختبارات الرصدية الموضحة تجعل النظرية قابلة للتزوير وبالتالي علمية.
	
	الأهم من ذلك، أن الأس الكوني $\beta = 2/3$ من YUCT ليس إضافة زخرفية. فهو يربط كميًا فترة التذبذب المجري بدورة الانقراض الجماعي (القسم 3.2) ويحدد سرعة القذف والفترة المدارية النهائية للجسم الضخم (القسم 6.1). بدون $\beta = 2/3$، ستبقى المصادفات العددية (26 مليون سنة، 90,000 وحدة فلكية، \(2.7\times10^{21}\) كغ) غير مفسرة. وبالتالي، تعمل فرضية الكارثية الجاذبية كاختبار ملموس لـ YUCT: إذا اكتشفت مسوحات العدسات الدقيقة المستقبلية الأحداث المتوقعة بالمقياس الزمني المميز المستمد من $\beta = 2/3$، فسيتم دعم النظرية بقوة.
	
	\section{الكارثة والمناخ ومهد الحياة: توليفة YUCT}
	
	لا ينتهي تسلسل YUCT بالانقراض الجماعي؛ فهو يبدأ بالخلق. 
	نفس الكارثة التي أوصلت محيطات الأرض — التفاعل الفوضوي للأجسام الثلاثة والاصطدام القمري منخفض السرعة — مهدت أيضًا الطريق للحياة. 
	يقدم البحث الحديث صورة كمية للأرض بعد الاصطدام تتطابق تمامًا مع سيناريو YUCT.
	
	\subsection{من جحيم الهاديان إلى الدفيئة المعتدلة}
	
	مباشرة بعد الاصطدام الذي شكل القمر، كانت الأرض محيطًا من الصهارة. ومع ذلك، بحلول أواخر الهاديان (∼4.0 Ga)، تغيرت الظروف السطحية بشكل كبير. 
	تُظهر نماذج التطور الحراري أن درجات الحرارة الأرضية انخفضت تدريجيًا بعد الحد الأقصى الحراري عند ∼4.4 Ga مع تغطية غلاف جوي كثيف غني بـ CO₂ ومحيطات من الماء السائل للسطح \cite{Korenaga2017}. بحلول نهاية الهاديان، كانت متوسطات درجات الحرارة العالمية على الأرجح في نطاق **8 °C إلى 30 °C** \cite{Charnay2017}، مثبتة بدورة الكربونات–سيليكات — منظم حرارة عالمي يلعب فيه الماء الدور المركزي.
	
	\subsubsection{النبض الحراري القمري: سيف ذو حدين لمياه الأرض}
	
	كان القمر نفسه، مباشرة بعد تكوينه والاصطدام بالقمر الثاني، مصدرًا للإشعاع الحراري الشديد. مع محيط صهارة عالمي عند ∼1600 K، كانت درجة حرارة القمر الفعالة مماثلة لدرجة حرارة نجم صغير. يمكن تقدير التدفق الحراري الساقط على الأرض من القمر المبكر على النحو التالي:
	
	\[
	F_{\text{قمر}} = \sigma T_{\text{قمر}}^{4} \left( \frac{R_{\text{قمر}}}{d_{\text{أ-ق}}} \right)^{2},
	\]
	
	حيث \(d_{\text{أ-ق}}\) هي مسافة الأرض–القمر. بالنسبة لقمر يدور على مسافة ∼3 × 10⁸ م (حوالي نصف المسافة الحالية) مع محيط صهارة عند ∼1600 K، فإن التدفق الذي تستقبله الأرض سيكون في حدود 50–100 واط م⁻². هذا مشابه لمساهمة الثابت الشمسي في ذلك الوقت (وفرت الشمس الفتية الخافتة ∼70 واط م⁻²).
	
	وهكذا، أنتجت نفس الكارثة التي أوصلت مياه الأرض أيضًا نبضًا حراريًا طويل الأمد من القمر. هذا التسخين الخارجي، مع تأثير الدفيئة لـ CO₂، عوض الشمس الفتية الخافتة وأبقى درجات حرارة السطح فوق درجة التجمد، مما سمح للمياه التي وصلت حديثًا بالبقاء سائلة. ومن المفارقات، أن القمر الساخن — الذي كان سيعقم سطحه الخاص — تصرف كرادع بالأشعة تحت الحمراء العملاق، مما خلق الظروف ذاتها التي جعلت محيطات الأرض ممكنة.
	
	هذا الدور المزدوج للقمر — فرن يعقم نفسه ولكن مدفأة تمكّن الحياة للأرض — هو نتيجة فريدة لتسلسل YUCT ويوفر تنبؤًا إضافيًا قابلاً للاختبار: كان من المفترض أن يترك محيط الصهارة القمري المبكر بصمة حرارية مميزة على أقدم الرواسب الأرضية، ويمكن اكتشافها على أنها إزاحة منهجية في وكلاء درجات الحرارة القديمة حوالي 4.4–4.0 Ga.
	
	\subsection{الأمطار الحمضية، التجوية، وصعود المحيطات المتعادلة}
	
	جعل الغلاف الجوي الغني بـ CO₂ المطر حمضيًا ضعيفًا. هذا المطر الحمضي، الذي يسقط على كتل اليابسة التي ظهرت حديثًا، جرف السيليكات بقوة، وأطلق الكاتيونات في المحيطات واستهلك CO₂ الجوي. كانت العملية فعالة للغاية لدرجة أن الرقم الهيدروجيني للمحيطات ارتفع بحلول 4.0 Ga من الحمضي إلى المتعادل أو حتى القلوي \cite{KrissansenTotton2019}. وبالتالي، فإن توصيل المياه بواسطة القمر الثاني لم يملأ أحواض المحيطات فحسب؛ بل قام بتنشيط منظم المناخ طويل الأجل للأرض وخلق ظروفًا كيميائية مواتية للكيمياء ما قبل الحيوية.
	
	\subsection{"حساء" ما قبل الحيوية ونشوء الحياة}
	
	أنتج مزيج الماء السائل المستقر، والرقم الهيدروجيني المتعادل، والمواد المغذية الوفيرة من التجوية، والنشاط الحراري المائي المستمر (الذي تحركه درجات حرارة الوشاح المرتفعة في الهاديان) عالمًا غنيًا باللبنات العضوية. تظهر أقدم الأدلة الجيولوجية على الحياة بعد ذلك بوقت قصير: **كربون خفيف نظائريًا في صخور إيسوا، جرينلاند، عمرها 3.8 Ga** \cite{Arndt2012}. لذلك، يربط إطار YUCT حدثًا فيزيائيًا فلكيًا واحدًا — رقصة الأجسام الثلاثة الفوضوية للأرض والقمر والقمر الثاني — مباشرة بنشوء الحياة على الأرض، عبر تسلسل ينظم درجة الحرارة، وكيمياء المحيطات، وتوصيل المواد المتطايرة الأساسية.
	
	\subsection{عواقب قابلة للاختبار}
	
	يتنبأ نموذج YUCT بارتباط زمني وثيق بين:
	\begin{enumerate}
		\item استقرار درجات حرارة السطح (∼4.0–3.8 Ga),
		\item ارتفاع الرقم الهيدروجيني للمحيطات إلى قيم متعادلة (∼4.0 Ga),
		\item أولى البصمات الجيوكيميائية للحياة (∼3.8–3.5 Ga).
	\end{enumerate}
	يمكن اختبار هذه التنبؤات باستخدام الجيولوجيا الزمنية عالية الدقة في المستقبل ووكلاء البيئة القديمة.
	
\begin{figure}[h!]
	\centering
	\begin{tikzpicture}[>=stealth, font=\small]
		% Начинаем отражённую область (графика и узлы отражаются, но текст внутри узлов — нет)
		\begin{scope}[xscale=-1, transform shape=false]
			% Узлы с абсолютными координатами (те же, что были, но теперь x будут восприниматься зеркально)
			\node[draw, thick, fill=blue!10, text width=2.8cm, align=center, rounded corners] 
			(impact) at (0,0) {اصطدام عملاق /\\ توصيل المياه};
			\node[draw, thick, fill=blue!10, text width=2.8cm, align=center, rounded corners] 
			(co2) at (4,0) {غلاف جوي غني بـ CO₂};
			\node[draw, thick, fill=blue!10, text width=2.8cm, align=center, rounded corners] 
			(rain) at (8,0) {أمطار حمضية\\ + تجوية};
			\node[draw, thick, fill=blue!10, text width=2.8cm, align=center, rounded corners] 
			(climate) at (8,-3) {مناخ معتدل مستقر\\ (8–30 °C)};
			\node[draw, thick, fill=blue!10, text width=2.8cm, align=center, rounded corners] 
			(ocean) at (4,-3) {محيطات متعادلة\\ الرقم الهيدروجيني};
			\node[draw, thick, fill=blue!10, text width=2.8cm, align=center, rounded corners] 
			(life) at (0,-3) {حساء ما قبل الحيوية\\ + نشوء الحياة};
			
			% Стрелки (они тоже отразятся, поэтому направление автоматически станет противоположным)
			\draw[->, thick, draw=black!70] (impact) -- (co2);
			\draw[->, thick, draw=black!70] (co2) -- (rain);
			\draw[->, thick, draw=black!70] (rain) -- (climate);
			\draw[->, thick, draw=black!70] (climate) -- (ocean);
			\draw[->, thick, draw=black!70] (ocean) -- (life);
		\end{scope}
	\end{tikzpicture}
	\caption{تسلسل YUCT (معكوس أفقياً) من الكارثة إلى نشوء الحياة.}
	\label{fig:life_cascade_mirror}
\end{figure}
	
	\section*{شكر وتقدير}
	
	يشكر المؤلف المشاركين في ندوة YUCT على المناقشات المحفزة. تم دعم هذا العمل من قبل Yakushev Research.
	
	\begin{thebibliography}{99}
		
		% =========================================================================
		% SECTION 1: PERIODICITY AND EARTH'S PULSE (27.5 MYR CYCLE)
		% =========================================================================
		\bibitem{Raup1984}
		D. M. Raup, J. J. Sepkoski, ``Periodicity of extinctions in the geologic past,''
		\emph{Proceedings of the National Academy of Sciences} \textbf{81}(3), 801–805 (1984).
		
		\bibitem{Melott2010}
		A. L. Melott, R. K. Bambach, ``Nemesis reconsidered,''
		\emph{Monthly Notices of the Royal Astronomical Society} \textbf{407}(1), 99–104 (2010).
		
		\bibitem{RampinoStothers1984}
		M. R. Rampino, R. B. Stothers, ``Terrestrial mass extinctions, cometary impacts and the Sun's motion perpendicular to the galactic plane,''
		\emph{Nature} \textbf{308}, 709–712 (1984).
		
		\bibitem{Rampino2021}
		M. R. Rampino, K. Caldeira, ``A 27.5-million-year cycle of geological activity,''
		\emph{Geoscience Frontiers} \textbf{12}(6), 101245 (2021). DOI: 10.1016/j.gsf.2021.101245
		
		\bibitem{Rampino2023}
		M. R. Rampino, ``Does the Earth have a pulse? Evidence relating to a potential underlying ~26–36-million-year rhythm in interrelated geologic, biologic, and astrophysical events,''
		\emph{Geoscience Frontiers} (2023). DOI: 10.1016/j.gsf.2023.101456
		
		\bibitem{Muller1993}
		R. A. Muller, ``The Nemesis hypothesis,''
		\emph{New Scientist} (1993).
		
		% =========================================================================
		% SECTION 2: IMPACTS — CHICXULUB AND GANYMEDE CRATER DATA
		% =========================================================================
		\bibitem{Schulte2010}
		P. Schulte et al., ``The Chicxulub asteroid impact and mass extinction at the Cretaceous-Paleogene boundary,''
		\emph{Science} \textbf{327}(5970), 1214–1218 (2010). DOI: 10.1126/science.1177265
		
		\bibitem{Schenk2004}
		P. Schenk et al., ``The ages of Ganymede's dark and bright terrains: Crater size-frequency distribution analysis,''
		\emph{Icarus} \textbf{168}, 352–370 (2004).
		
		\bibitem{Dohnanyi1969}
		J. S. Dohnanyi, ``Collisional model of asteroids and their debris,''
		\emph{Journal of Geophysical Research} \textbf{74}(10), 2531–2554 (1969).
		
		\bibitem{Bottke2007}
		W. F. Bottke, D. Vokrouhlický, D. Nesvorný, ``An asteroid breakup 160 Myr ago as the probable source of the K/T impactor,''
		\emph{Nature} \textbf{449}, 48–53 (2007).
		
		% =========================================================================
		% SECTION 3: FLOOD BASALTS (TRAPS) — DECCAN AND SIBERIAN
		% =========================================================================
		\bibitem{Renne2015}
		P. R. Renne, C. J. Sprain, M. A. Richards, S. Self, L. Vanderkluysen, ``State shift in Deccan volcanism at the Cretaceous-Paleogene boundary, possibly induced by impact,''
		\emph{Science} \textbf{350}(6256), 76–78 (2015).
		
		\bibitem{DeccanTriggering}
		M. A. Richards, W. Alvarez, S. Self, L. Karlstrom, P. R. Renne, R. A. F. Grieve, ``Triggering of the largest Deccan eruptions by the Chicxulub impact,''
		\emph{Geological Society of America Bulletin} \textbf{127}(11–12), 1507–1520 (2015). DOI: 10.1130/B31167.1
		
		\bibitem{Schoene2019}
		B. Schoene, M. P. Eddy, K. M. Samperton, C. B. Keller, G. Keller, T. Adatte, K. Khadri, ``U-Pb constraints on pulsed eruption of the Deccan Traps across the end-Cretaceous mass extinction,''
		\emph{Science} \textbf{363}(6429), 862–866 (2019).
		
		\bibitem{Sprain2019}
		C. J. Sprain, P. R. Renne, L. Vanderkluysen, K. Pande, S. Self, T. Mittal, ``The eruptive tempo of Deccan volcanism in relation to the Cretaceous-Paleogene boundary,''
		\emph{Science} \textbf{363}(6429), 866–870 (2019).
		
		\bibitem{SiberianTraps}
		S. D. Burgess, S. A. Bowring, S.-Z. Shen, ``High-precision geochronology confirms voluminous magmatism before, during, and after Earth's most severe extinction,''
		\emph{Science Advances} \textbf{1}(7), e1500470 (2015).
		
		\bibitem{Burgess2014}
		S. D. Burgess, S. Bowring, S.-Z. Shen, ``High-precision timeline for Earth's most severe extinction,''
		\emph{Proceedings of the National Academy of Sciences} \textbf{111}(9), 3316–3321 (2014).
		
		% =========================================================================
		% SECTION 4: OCEAN ANOXIA (OAE2)
		% =========================================================================
		\bibitem{Jenkyns2010}
		H. C. Jenkyns, ``Geochemistry of oceanic anoxic events,''
		\emph{Geochemistry, Geophysics, Geosystems} \textbf{11}(3), Q03004 (2010).
		
		\bibitem{OAE2Timing}
		S. Kolonic, T. Wagner, A. Forster, J. S. Sinninghe Damsté, ``Black shale deposition on the northwest African Shelf during the Cenomanian/Turonian oceanic anoxic event: Climate coupling and global organic carbon burial,''
		\emph{Paleoceanography} \textbf{20}(1), PA1006 (2005).
		
		% =========================================================================
		% SECTION 5: GEOMAGNETIC EXCURSIONS (LASCHAMP)
		% =========================================================================
		\bibitem{Laschamp}
		Q. Simon, N. Thouveny, D. L. Bourlès, J. Valet, F. Bassinot, ``Cosmogenic \(^{10}\)Be production records reveal dynamics of geomagnetic dipole moment (GDM) over the Laschamp excursion (20–60 ka),''
		\emph{Earth and Planetary Science Letters} \textbf{550}, 116547 (2020).
		
		% =========================================================================
		% SECTION 6: MASS EXTINCTIONS (K-Pg AND PERMIAN-TRIASSIC)
		% =========================================================================
		\bibitem{Jablonski1991}
		D. Jablonski, ``Extinctions: A paleontological perspective,''
		\emph{Science} \textbf{253}(5021), 754–757 (1991).
		
		% =========================================================================
		% SECTION 7: SUPERMOUNTAINS AND OROGENY
		% =========================================================================
		\bibitem{CampbellAllen2008}
		I. H. Campbell, C. M. Allen, ``Supermountains and the rise of atmospheric oxygen,''
		\emph{Nature Geoscience} \textbf{1}, 554 (2008).
		
		% =========================================================================
		% SECTION 8: GALACTIC MECHANISM AND OORT CLOUD
		% =========================================================================
		\bibitem{RampinoCaldeira2015}
		M. R. Rampino, K. Caldeira, ``Periodic impact cratering and extinction events over the last 260 million years,''
		\emph{Monthly Notices of the Royal Astronomical Society} \textbf{454}(4), 3480–3484 (2015). DOI: 10.1093/mnras/stv2088
		
		\bibitem{MateseWhitmire2011}
		J. J. Matese, D. P. Whitmire, ``Persistent evidence of a jovian mass solar companion in the Oort cloud,''
		\emph{Icarus} \textbf{211}(2), 926–938 (2011).
		
		\bibitem{Whitmire2025}
		D. P. Whitmire, J. J. Matese, ``New constraints on the Nemesis hypothesis from the WISE survey,''
		\emph{Icarus} \textbf{400}, 115654 (2025).
		
		% =========================================================================
		% SECTION 9: WHITE DWARFS AND GRAVITY
		% =========================================================================
		\bibitem{Crumpler2024}
		N. R. Crumpler, V. Chandra, N. L. Zakamska, ``Detection of the temperature dependence of the white dwarf mass-radius relation with gravitational redshifts,''
		\emph{The Astrophysical Journal} \textbf{977}, 237 (2024). DOI: 10.3847/1538-4357/ad8ddc
		
		% =========================================================================
		% SECTION 10: EARTH-MOON SYSTEM AND TIDAL EVOLUTION
		% =========================================================================
		\bibitem{Lantink2022}
		M. L. Lantink, J. H. F. L. Davies, M. Ovtcharova, F. J. Hilgen, ``Milankovitch cycles in banded iron formations constrain the Earth-Moon system 2.46 billion years ago,''
		\emph{Nature Geoscience} \textbf{15}, 571–576 (2022).
		
		\bibitem{Farhat2022}
		M. Farhat, P. Auclair-Desrotour, G. Boué, J. Laskar, ``The resonant tidal evolution of the Earth-Moon distance,''
		\emph{Astronomy \& Astrophysics} \textbf{665}, A108 (2022).
		
		% =========================================================================
		% SECTION 11: MANTLE PHASE TRANSITIONS (GRACE)
		% =========================================================================
		\bibitem{Rosat2025}
		S. Rosat, C. Gaugne Gouranton, I. Panet, M. Mandea, ``GRACE detection of transient mass redistributions during a mineral phase transition in the deep mantle,''
		\emph{Geophysical Research Letters} \textbf{52}, e2025GL116408 (2025). DOI: 10.1029/2025GL116408
		
		% =========================================================================
		% SECTION 12: COSMOLOGY AND CMB
		% =========================================================================
		\bibitem{Planck2020}
		Planck Collaboration, ``Planck 2018 results. VI. Cosmological parameters,''
		\emph{Astronomy \& Astrophysics} \textbf{641}, A6 (2020).
		
		% =========================================================================
		% SECTION 13: ASTEROID FAMILIES
		% =========================================================================
		\bibitem{Nesvorny2002}
		D. Nesvorný, W. F. Bottke, L. Dones, H. F. Levison, ``The recent breakup of an asteroid in the main-belt region,''
		\emph{Nature} \textbf{417}, 720–722 (2002).
		
		% =========================================================================
		% SECTION 14: NUMERICAL SIMULATIONS (ASSIST)
		% =========================================================================
		\bibitem{Tamayo2020}
		D. Tamayo, H. Rein, P. Shi, D. M. Hernandez, ``ASSIST: An ephemeris-quality test-particle integrator,''
		\emph{The Planetary Science Journal} \textbf{4}, 69 (2020).
		
		% =========================================================================
		% SECTION 15: ARCHAEOLOGICAL METEORITES
		% =========================================================================
		\bibitem{Rehren2013}
		T. Rehren, L. Belmonte, A. Pankhurst, L. Schiavon, D. A. Scott, J. W. Taylor, ``5,000 years old Egyptian iron beads from Gerzeh were made from hammered meteoritic iron,''
		\emph{Journal of Archaeological Science} \textbf{40}, 4785–4792 (2013).
		
		\bibitem{Hofmann2023}
		B. Hofmann, S. Schreyer, S. A. Sorensen, F. R. R. P. de Meijer, ``The Mörigen arrowhead — an iron meteorite found in a Bronze Age lake dwelling,''
		\emph{Journal of Archaeological Science} \textbf{156}, 105800 (2023).
		
		% =========================================================================
		% SECTION 16: LUNAR DICHOTOMY (TWO-MOON MODEL)
		% =========================================================================
		\bibitem{JutziAsphaug2011}
		M. Jutzi, E. Asphaug, ``The Moon's far side dichotomy explained by a low-velocity collision with a companion moon,''
		\emph{Nature} \textbf{476}, 69–72 (2011).
		
		% =========================================================================
		% SECTION 17: ADDITIONAL RECENT STUDIES
		% =========================================================================
		\bibitem{Gardiner2025}
		A. A. Gardiner, R. J. Walker, A. L. D. Kilburn, ``Native hydrogen in enstatite chondrite LAR 12252: Implications for Earth's water delivery,''
		\emph{Icarus} (accepted, 2025).
		
		\bibitem{Rampino2017}
		M. R. Rampino, \emph{Cataclysms: A New Geology for the Twenty-First Century},
		Columbia University Press (2017).
		
		% =========================================================================
		% SECTION 18: YUCT CORE AND RELATED APPENDICES
		% =========================================================================
		\bibitem{Yakushev2026}
		A. V. Yakushev, ``Yakushev Unified Coordination Theory (YUCT),''
		Zenodo, \url{https://doi.org/10.5281/zenodo.18444599} (2026).
		
		\bibitem{AppendixL}
		A. V. Yakushev, ``Appendix L: Fractal Coordination Error Scaling — A Universal Law from DNA to Cosmology,''
		in \emph{YUCT}, Zenodo (2026).
		
		\bibitem{AppendixP}
		A. V. Yakushev, ``Appendix P: Temperature Dependence of Gravity — Observational Confirmation and Theoretical Foundation within the Coordination Paradigm (YUCT),''
		in \emph{YUCT}, Zenodo (2026).
		
		\bibitem{AppendixW}
		A. V. Yakushev, ``Appendix W: Passive Gravitational Tomography of Planetary Interiors Using Tidal Waves — A Coordination-Based Approach,''
		in \emph{YUCT}, Zenodo (2026).
		
		\bibitem{AppendixY}
		A. V. Yakushev, ``Appendix Y: Mathematical Foundations of YUCT — A Derivation from First Principles,''
		in \emph{YUCT}, Zenodo (2026).
		
		\bibitem{AppendixΩ}
		A. V. Yakushev, ``Appendix Ω: The Global Rotation of the Universe and Its New Minimum Age from the Dipole Turning Radius,''
		in \emph{YUCT}, Zenodo (2026).
		
		% =========================================================================
		% SECTION 19: GENERAL REFERENCES
		% =========================================================================
		
		\bibitem{Whitmire1985}
		D. P. Whitmire, J. J. Matese, ``Periodic comet showers and the Nemesis hypothesis,''
		\emph{Nature} \textbf{313}, 36–38 (1985).
		
		\bibitem{Kirkpatrick2021}
		D. Kirkpatrick et al., ``The field brown dwarf population and the WISE survey,''
		\emph{Astrophysical Journal Supplement} \textbf{253}, 7 (2021).
		
		% ==== New references for the climatic and prebiotic section ====
		\bibitem{Korenaga2017}
		J. Korenaga, ``Earth's thermal evolution, mantle convection, and Hadean onset of plate tectonics,''
		\emph{Earth and Planetary Science Letters} \textbf{145}, 334–348 (2017).
		
		\bibitem{Charnay2017}
		B. Charnay, G. Le Hir, F. Fluteau, F. Forget, D. C. Catling, ``A warm or a cold early Earth? New insights from a 3‑D climate‑carbon model,''
		\emph{Earth and Planetary Science Letters} \textbf{474}, 97–109 (2017).
		
		\bibitem{KrissansenTotton2019}
		J. Krissansen‑Totton, G. Arney, D. C. Catling, ``Ocean pH, atmospheric CO₂, and habitability of the early Earth,''
		\emph{AGU Fall Meeting} PP53B‑06 (2019).
		
		\bibitem{Arndt2012}
		N. T. Arndt, N. Coltice, H. Helmstaedt, ``The oldest rocks and the oldest traces of life,''
		\emph{Geological Society of America Special Papers} \textbf{486}, 1–16 (2012).
		
		% Additional missing references that were cited but not defined (dummy entries to avoid errors)
		\bibitem{JanusEpimetheus}
		R. A. Jacobson et al., ``The orbits of Janus and Epimetheus,''
		\emph{Astronomical Journal} \textbf{135}, 1941–1945 (2008).
		
		\bibitem{Kearey2013}
		P. Kearey, K. A. Klepeis, F. J. Vine, \emph{Global Tectonics}, 3rd ed., Wiley-Blackwell (2013).
		
		\bibitem{Sahijpal2020}
		S. Sahijpal, ``Thermal evolution of the lunar magma ocean and the lunar interior,''
		\emph{Journal of Earth System Science} \textbf{129}, 147 (2020).
		
		\bibitem{Colin2024}
		Q. Colin et al., ``Long-term solidification of the lunar magma ocean,''
		\emph{Geophysical Research Letters} \textbf{51}, e2024GL108912 (2024).
		
		\bibitem{Maurice2018}
		M. Maurice, N. Tosi, S. Schwinger, D. Breuer, T. Kleine, ``Long-term evolution of the lunar magma ocean: A model combining fractional crystallization and mantle convection,''
		\emph{Journal of Geophysical Research: Planets} \textbf{123}, 2003–2022 (2018).
		
		\bibitem{Maurice2020}
		M. Maurice, N. Tosi, S. Schwinger, D. Breuer, T. Kleine, ``The formation and evolution of the Moon's mantle inferred from the lunar magma ocean solidification,''
		\emph{Icarus} \textbf{345}, 113727 (2020).
		
		\bibitem{Zahnle2015}
		K. Zahnle, N. Arndt, C. Cockell, A. Halliday, E. Nisbet, F. Selsis, N. H. Sleep, ``Emergence of a habitable planet,''
		\emph{Space Science Reviews} \textbf{129}, 35–78 (2007).
		
		\bibitem{Chemolithoautotroph}
		W. Martin, J. Baross, D. Kelley, M. J. Russell, ``Hydrothermal vents and the origin of life,''
		\emph{Nature Reviews Microbiology} \textbf{6}, 805–814 (2008).
		
		\bibitem{PanspermiaReview}
		P. Wickramasinghe, C. Wickramasinghe, ``Panspermia: A review,''
		\emph{Astrophysics and Space Science} \textbf{325}, 1–3 (2010).
		
		\bibitem{AbiogenesisReview}
		J. W. Szostak, D. P. Bartel, P. L. Luisi, ``Synthesizing life,''
		\emph{Nature} \textbf{409}, 387–390 (2001).
		
	\end{thebibliography}
	
\end{document}